\documentclass[a4paper,11pt,oneside,openany]{book}
\usepackage{exam_style}

\title{ECE181: Introduction to Python --- Examination Companion \& Solved Question Papers}
\author{LPU Engineering Academic Faculty}
\date{Edition 2026--2027}

\begin{document}

\frontmatter
\begin{titlepage}
\pagecolor{LPUNavy}
\color{white}
\vspace*{1.5cm}

\begin{center}
    {\Huge \bfseries \scshape ECE181: Introduction to Python}\\[5mm]
    {\LARGE \bfseries The Definitive University Examination Companion \& Solved Model Papers}\\[8mm]
    {\color{LPUOrange}\rule{0.8\textwidth}{2pt}}\\[10mm]

    \begin{tcolorbox}[enhanced, width=0.88\textwidth, colback=white, colframe=LPUOrange, boxrule=1.5pt, arc=3mm, halign=center]
        \color{DarkSlate}
        \vspace{2mm}
        {\Large\bfseries Complete Assessment Suite (Sets A \& B for CA-1, Mid-Term, CA-2, End-Term)}\\[3mm]
        {\large\bfseries Featuring Step-by-Step Code Implementations \& Official Marking Rubrics}\\[3mm]
        \begin{tabular}{cc}
            \color{LPUNavy}\textbf{Full CA-1 Papers (Units 1 \& 2)} & \quad \color{LPUNavy}\textbf{Full Mid-Term Papers (Units 1, 2, 3)} \\
            \color{LPUNavy}\textbf{Full CA-2 Papers (Units 4 \& 5)} & \quad \color{LPUNavy}\textbf{Full End-Term Papers (Units 1 to 6)} \\
        \end{tabular}
        \vspace{2mm}
    \end{tcolorbox}

    \vfill

    {\large \bfseries Prepared for Lovely Professional University (LPU) Engineering Students}\\[2mm]
    {\small Course Outcomes CO1--CO6 $\bullet$ Bloom's Taxonomy L1--L4 $\bullet$ Exact Marking Rubrics}\\[10mm]

    \begin{tcolorbox}[width=0.6\textwidth, colback=LPUNavy!80!black, colframe=white, boxrule=0.5pt, halign=center]
        \color{white}
        \textbf{LPU Engineering Academic Series}\\[1mm]
        {\footnotesize Academic Session 2026--2027}
    \end{tcolorbox}
    \vspace*{1cm}
\end{center}
\end{titlepage}
\nopagecolor
\newpage

\chapter*{Official Examination Architecture \& Blueprint}
\addcontentsline{toc}{chapter}{Official Examination Architecture \& Blueprint}

This Examination Companion provides verified model question papers and complete step-by-step solutions designed strictly to mirror the Lovely Professional University (LPU) examination structure for \textbf{ECE181: Introduction to Python}:

\begin{center}
\begin{tabularx}{\textwidth}{l p{3.2cm} p{2.8cm} c X}
\toprule
\textbf{Assessment} & \textbf{Syllabus Scope} & \textbf{Format} & \textbf{Marks} & \textbf{Instructions \& Rules} \\
\midrule
\textbf{CA-1} & Unit 1 \& Unit 2 & 6 Subjective Qs & 30 & Attempt any 5 out of 6 questions ($5 \times 6\text{M} = 30\text{M}$). Time: 60 mins. \\
\addlinespace
\textbf{Mid-Term} & Unit 1, Unit 2, Unit 3 & 30 MCQs & 30 & 10 MCQs per unit. Negative marking $-0.25$ marks for wrong answers. Time: 90 mins. \\
\addlinespace
\textbf{CA-2} & Unit 4 \& Unit 5 & 6 Subjective Qs & 30 & Attempt any 5 out of 6 questions ($5 \times 6\text{M} = 30\text{M}$). Time: 60 mins. \\
\addlinespace
\textbf{End-Term} & Units 1 to 6 & \textbf{Part A}: 30 MCQs \newline \textbf{Part B}: 5 Subj Qs & 70 & \textbf{Part A}: 30 MCQs from Units 4, 5, 6 (10 each). \newline \textbf{Part B}: 5 Qs of 10M, attempt any 4 ($4 \times 10 = 40\text{M}$). Time: 3 Hours. \\
\bottomrule
\end{tabularx}
\end{center}

\newpage

\chapter*{Core Course Outcomes \& Bloom's Taxonomy Guide}
\addcontentsline{toc}{chapter}{Core Course Outcomes \& Bloom's Taxonomy Guide}

\begin{itemize}[leftmargin=1.5em, itemsep=4pt]
    \item \textbf{CO1}: Understand procedural Python syntax, variable binding, numeric primitives, and operator precedence rules.
    \item \textbf{CO2}: Construct robust control-flow algorithms utilizing conditional branching, structured loops, and recursive function designs.
    \item \textbf{CO3}: Select, manipulate, and optimize composite linear and non-linear data structures (strings, lists, tuples, and dictionaries).
    \item \textbf{CO4}: Leverage NumPy N-dimensional arrays, vectorization, and multi-dimensional broadcasting for high-performance numerical computing.
    \item \textbf{CO5}: Perform structured tabular data processing and aggregation using Pandas, and produce publication-grade statistical charts with Matplotlib.
    \item \textbf{CO6}: Implement integrated Python pipelines connecting digital image processing (OpenCV), interactive GUI loops (Pygame), and automated web scraping.
\end{itemize}

\vspace{5mm}
\noindent
\textbf{Bloom's Taxonomy Levels Utilized in Question Marking:}
\begin{enumerate}[label=\textbf{L\arabic*:}, leftmargin=1.5em, itemsep=2pt]
    \item \textbf{Remember:} Stating definitions, keywords, method signatures, and default behaviors.
    \item \textbf{Understand:} Explaining control flow, tracing variable scope (LEGB), and predicting output of code snippets.
    \item \textbf{Apply:} Implementing algorithmic logic, correcting syntax traps, and computing transformations.
    \item \textbf{Analyze:} Formulating time and space complexity $\mathcal{O}(n)$, diagnosing memory mutations, and optimizing vectorized operations.
\end{enumerate}

\newpage


\tableofcontents
\newpage

\mainmatter

% ==============================================================================
% SECTION 1: CONTINUOUS ASSESSMENT 1 (CA-1) --- UNITS 1 & 2
% ==============================================================================
\chapter{Continuous Assessment 1 (CA-1) Suite}
\label{chap:ca1}
\noindent
Continuous Assessment 1 evaluates students on \textbf{Unit 1 (Introduction, Variables, Expressions)} and \textbf{Unit 2 (Conditionals, Loops, Functions \& Recursion)}. Each paper contains \textbf{6 questions} carrying \textbf{6 marks each}, and students attempt \textbf{5 questions} (Total = \textbf{30 Marks}, Time Allowed = \textbf{60 Minutes}).

\newpage
\section{Continuous Assessment 1 (CA-1) --- Set A}
\label{sec:ca1_set_a}

\LPUPaperHeader{Continuous Assessment (CA-1)}{A}{60 Minutes}{30 Marks}{CO1 (Foundations \& Variables), CO2 (Control Flow \& Functions)}

\begin{center}
\sffamily\textbf{\color{AlertRed}Note: Attempt any five out of the six questions. Each question carries 6 marks.}
\end{center}

\questionitem{1}{
    Explain the difference between dynamic typing in Python and static typing in C/Java. Trace how the Python Virtual Machine (PVM) executes \texttt{x = 5} followed by \texttt{x = "ECE181"}, detailing heap object allocation and reference counting.
}{6}{CO1}{L2 Understand}

\questionitem{2}{
    Evaluate the following arithmetic and bitwise expressions step by step, applying exact operator precedence and associativity rules:
    \begin{enumerate*}[label=(\alph*)]
        \item \texttt{2 ** 3 ** 2} \quad
        \item \texttt{-17 // 5} \quad
        \item \texttt{17 \% -5} \quad
        \item \texttt{5 << 2 + 1}
    \end{enumerate*}
}{6}{CO1}{L3 Apply}

\questionitem{3}{
    Write a Python program that accepts a student's numerical marks $[0, 100]$ and outputs their letter grade using an optimized \texttt{if-elif-else} construct. Include defensive input validation that raises a \texttt{ValueError} if marks are outside the valid domain or of an invalid type.
}{6}{CO2}{L3 Apply}

\questionitem{4}{
    Differentiate between \texttt{break} and \texttt{continue} statements. Write a Python snippet that prints all prime numbers between $10$ and $50$ using a \texttt{for-else} loop construct, explaining the exact condition under which the \texttt{else} block triggers.
}{6}{CO2}{L3 Apply}

\questionitem{5}{
    Explain the \textbf{Mutable Default Argument} trap in Python function definitions. Provide a code example demonstrating how \texttt{def append\_to(element, target=[])} produces unexpected cross-call state retention, and write the corrected canonical idiom.
}{6}{CO2}{L4 Analyze}

\questionitem{6}{
    Implement a recursive function \texttt{power(base, exp)} that computes $\text{base}^{\text{exp}}$ in $\mathcal{O}(\log \text{exp})$ time using the method of exponentiation by squaring. Write the recurrence relation and draw the call stack for \texttt{power(3, 5)}.
}{6}{CO2}{L4 Analyze}

\newpage
\subsection*{Solutions and Marking Rubrics --- CA-1 Set A}

\begin{solutionbox}[Solution to Question 1]
In statically typed languages (C/Java), types are bound to \textbf{variable containers} at compile-time: \texttt{int x = 5;} allocates 4 bytes of stack memory. In Python, variables are untyped \textbf{names (pointers/references)} bound to heap objects:
\begin{enumerate}[leftmargin=1.5em]
    \item \texttt{x = 5}: PVM allocates a \texttt{PyLongObject} on the heap storing value 5 with \texttt{ob\_refcnt = 1}. Name \texttt{x} in the local namespace dictionary maps to this address.
    \item \texttt{x = "ECE181"}: PVM allocates a new \texttt{PyUnicodeObject} on the heap. Name \texttt{x} is rebound to the string. The integer object's \texttt{ob\_refcnt} is decremented to 0, triggering automatic deallocation by Python's garbage collector.
\end{enumerate}
\begin{markingrubric}
$\bullet$ Distinction between variable-typed vs object-typed paradigms: \textbf{2 Marks}\\
$\bullet$ Heap object representation (\texttt{PyObject}, reference pointer): \textbf{2 Marks}\\
$\bullet$ Reference count modification and garbage collection mechanics: \textbf{2 Marks}
\end{markingrubric}
\end{solutionbox}

\begin{solutionbox}[Solution to Question 2]
\begin{enumerate}[label=(\alph*), leftmargin=1.5em]
    \item \texttt{2 ** 3 ** 2}: Exponentiation has right-to-left associativity: $2^{(3^2)} = 2^9 = \mathbf{512}$.
    \item \texttt{-17 // 5}: Floor division rounds down toward negative infinity: $\lfloor -17/5 \rfloor = \lfloor -3.4 \rfloor = \mathbf{-4}$.
    \item \texttt{17 \% -5}: Modulo identity $a \pmod b = a - b \times (a // b)$. Since $17 // -5 = -4$, $17 - (-5)(-4) = 17 - 20 = \mathbf{-3}$.
    \item \texttt{5 << 2 + 1}: Addition has precedence 5 while left shift has precedence 6. Evaluates as $5 \ll (2 + 1) = 5 \ll 3 = 5 \times 2^3 = \mathbf{40}$.
\end{enumerate}
\begin{markingrubric}
$\bullet$ Correct evaluation of (a) with right-associativity justification: \textbf{1.5 Marks}\\
$\bullet$ Correct evaluation of (b) and (c) with floor rounding rule: \textbf{3.0 Marks}\\
$\bullet$ Correct evaluation of (d) with operator precedence order: \textbf{1.5 Marks}
\end{markingrubric}
\end{solutionbox}

\begin{solutionbox}[Solution to Question 3]
\begin{pycode}
def assign\_grade(marks):\\
\_\_ESCTOKEN\_7\_\_ if not isinstance(marks, (int, float)):\\
\_\_ESCTOKEN\_7\_\_ \_\_ESCTOKEN\_7\_\_ raise TypeError("Marks must be an integer or float")\\
\_\_ESCTOKEN\_7\_\_ if not (0 <= marks <= 100):\\
\_\_ESCTOKEN\_7\_\_ \_\_ESCTOKEN\_7\_\_ raise ValueError("Marks must reside within closed interval [0, 100]")\\
~\\
\_\_ESCTOKEN\_7\_\_ if marks >= 90:\\
\_\_ESCTOKEN\_7\_\_ \_\_ESCTOKEN\_7\_\_ return "A"\\
\_\_ESCTOKEN\_7\_\_ elif marks >= 80:\\
\_\_ESCTOKEN\_7\_\_ \_\_ESCTOKEN\_7\_\_ return "B"\\
\_\_ESCTOKEN\_7\_\_ elif marks >= 70:\\
\_\_ESCTOKEN\_7\_\_ \_\_ESCTOKEN\_7\_\_ return "C"\\
\_\_ESCTOKEN\_7\_\_ elif marks >= 60:\\
\_\_ESCTOKEN\_7\_\_ \_\_ESCTOKEN\_7\_\_ return "D"\\
\_\_ESCTOKEN\_7\_\_ elif marks >= 40:\\
\_\_ESCTOKEN\_7\_\_ \_\_ESCTOKEN\_7\_\_ return "E"\\
\_\_ESCTOKEN\_7\_\_ else:\\
\_\_ESCTOKEN\_7\_\_ \_\_ESCTOKEN\_7\_\_ return "F"
\end{pycode}
\begin{markingrubric}
$\bullet$ Defensive type and boundary checking with explicit exceptions: \textbf{2 Marks}\\
$\bullet$ Non-overlapping, mutually exclusive \texttt{elif} ladder: \textbf{3 Marks}\\
$\bullet$ Correct return value syntax and cleanliness: \textbf{1 Mark}
\end{markingrubric}
\end{solutionbox}

\begin{solutionbox}[Solution to Question 4]
\textbf{Difference}: \texttt{break} terminates the loop completely, skipping any associated loop \texttt{else} block. \texttt{continue} skips only the remaining instructions of the current iteration, proceeding to the next loop step.
\begin{pycode}
for num in range(10, 51):\\
\_\_ESCTOKEN\_7\_\_ for d in range(2, int(num ** 0.5) + 1):\\
\_\_ESCTOKEN\_7\_\_ \_\_ESCTOKEN\_7\_\_ if num \% d == 0:\\
\_\_ESCTOKEN\_7\_\_ \_\_ESCTOKEN\_7\_\_ \_\_ESCTOKEN\_7\_\_ break\\
\_\_ESCTOKEN\_7\_\_ else:\\
\_\_ESCTOKEN\_7\_\_ \_\_ESCTOKEN\_7\_\_ print(num, end=" ")
\end{pycode}
\textbf{Condition for \texttt{else} Triggering}: The \texttt{else} clause fires if and only if the inner loop finishes all iterations across $2 \dots \lfloor \sqrt{num} \rfloor$ without executing a \texttt{break} statement, confirming $num$ is prime.
\begin{markingrubric}
$\bullet$ Precise conceptual difference between \texttt{break} and \texttt{continue}: \textbf{2 Marks}\\
$\bullet$ Correct prime checking logic utilizing \texttt{range(2, int(sqrt(num))+1)}: \textbf{2 Marks}\\
$\bullet$ Explanation and implementation of loop \texttt{else} trigger: \textbf{2 Marks}
\end{markingrubric}
\end{solutionbox}

\begin{solutionbox}[Solution to Question 5]
In Python, default parameter expressions are evaluated \textbf{once when the \texttt{def} statement is executed}, not each time the function is called.
\begin{pycode}
def append\_to(element, target=[]):\\
\_\_ESCTOKEN\_7\_\_ target.append(element)\\
\_\_ESCTOKEN\_7\_\_ return target\\
~\\
print(append\_to(1))  \# [1]\\
print(append\_to(2))  \# [1, 2] -> Unexpected shared mutation!
\end{pycode}
\textbf{Canonical Solution}:
\begin{pycode}
def append\_to(element, target=None):\\
\_\_ESCTOKEN\_7\_\_ if target is None:\\
\_\_ESCTOKEN\_7\_\_ \_\_ESCTOKEN\_7\_\_ target = []\\
\_\_ESCTOKEN\_7\_\_ target.append(element)\\
\_\_ESCTOKEN\_7\_\_ return target
\end{pycode}
\begin{markingrubric}
$\bullet$ Explanation of function object definition-time evaluation: \textbf{2 Marks}\\
$\bullet$ Demonstration of shared list mutation bug: \textbf{2 Marks}\\
$\bullet$ Correct implementation using \texttt{target=None} guard: \textbf{2 Marks}
\end{markingrubric}
\end{solutionbox}

\begin{solutionbox}[Solution to Question 6]
\begin{pycode}
def power(base, exp):\\
\_\_ESCTOKEN\_7\_\_ if exp == 0:\\
\_\_ESCTOKEN\_7\_\_ \_\_ESCTOKEN\_7\_\_ return 1.0\\
\_\_ESCTOKEN\_7\_\_ elif exp < 0:\\
\_\_ESCTOKEN\_7\_\_ \_\_ESCTOKEN\_7\_\_ return 1.0 / power(base, -exp)\\
\_\_ESCTOKEN\_7\_\_ half = power(base, exp // 2)\\
\_\_ESCTOKEN\_7\_\_ if exp \% 2 == 0:\\
\_\_ESCTOKEN\_7\_\_ \_\_ESCTOKEN\_7\_\_ return half * half\\
\_\_ESCTOKEN\_7\_\_ else:\\
\_\_ESCTOKEN\_7\_\_ \_\_ESCTOKEN\_7\_\_ return base * half * half
\end{pycode}
\textbf{Recurrence Relation}: $T(n) = T(n/2) + \mathcal{O}(1) \implies T(n) = \mathcal{O}(\log n)$.\\
\textbf{Call Stack for \texttt{power(3, 5)}}:
\begin{itemize}[leftmargin=1.5em]
    \item Frame 1: \texttt{power(3, 5)} $\to$ calls \texttt{power(3, 2)}
    \item Frame 2: \texttt{power(3, 2)} $\to$ calls \texttt{power(3, 1)}
    \item Frame 3: \texttt{power(3, 1)} $\to$ calls \texttt{power(3, 0)}
    \item Frame 4: \texttt{power(3, 0)} returns $1.0$ (base case)
    \item Unwinding: Frame 3 returns $3 \times 1 \times 1 = 3$; Frame 2 returns $3 \times 3 = 9$; Frame 1 returns $3 \times 9 \times 9 = 243.0$.
\end{itemize}
\begin{markingrubric}
$\bullet$ Logarithmic recursive implementation with negative exponent handling: \textbf{2 Marks}\\
$\bullet$ Recurrence relation derivation via Master Theorem: \textbf{2 Marks}\\
$\bullet$ Accurate call stack visualization and unwinding return values: \textbf{2 Marks}
\end{markingrubric}
\end{solutionbox}

\newpage

\section{Continuous Assessment 1 (CA-1) --- Set B}
\label{sec:ca1_set_b}

\LPUPaperHeader{Continuous Assessment (CA-1)}{B}{60 Minutes}{30 Marks}{CO1 (Variables \& Types), CO2 (Conditionals \& Functions)}

\begin{center}
\sffamily\textbf{\color{AlertRed}Note: Attempt any five out of the six questions. Each question carries 6 marks.}
\end{center}

\questionitem{1}{
    Explain the LEGB (Local, Enclosing, Global, Built-in) scope resolution rule. Provide code examples demonstrating the use and distinct behaviors of the \texttt{global} and \texttt{nonlocal} keywords.
}{6}{CO2}{L2 Understand}

\questionitem{2}{
    Discuss floating-point representation artifacts in Python IEEE 754 arithmetic (e.g., \texttt{0.1 + 0.2 != 0.3}). Explain how to correctly perform financial arithmetic using the \texttt{decimal} module and tolerance comparisons via \texttt{math.isclose()}.
}{6}{CO1}{L3 Apply}

\questionitem{3}{
    Write a Python function \texttt{collatz\_sequence(n)} that generates the sequence of numbers for starting integer $n$ until it reaches 1 ($n \to n/2$ if even, $n \to 3n + 1$ if odd). Compute and return the peak value reached and total stopping time.
}{6}{CO2}{L3 Apply}

\questionitem{4}{
    Analyze Python's short-circuit evaluation mechanism for boolean operators \texttt{and} and \texttt{or}. What is printed by \texttt{print([] or 0 or "Hello" and 5)}? Explain the evaluation sequence step-by-step.
}{6}{CO1}{L4 Analyze}

\questionitem{5}{
    Explain the function of \texttt{random.seed()}. Differentiate between pseudo-random numbers and cryptographically secure random numbers, contrasting \texttt{random.choice()} with \texttt{secrets.choice()}.
}{6}{CO2}{L2 Understand}

\questionitem{6}{
    Write a recursive function to solve the classical \textbf{Tower of Hanoi} puzzle for $n$ disks. Draw the recursion tree for $n = 3$, showing disk movements and deriving total move count $M(n) = 2^n - 1$.
}{6}{CO2}{L4 Analyze}

\newpage
\subsection*{Solutions and Marking Rubrics --- CA-1 Set B}

\begin{solutionbox}[Solution to Question 1]
When resolving an identifier, Python searches namespaces in strict order:
\begin{enumerate}[leftmargin=1.5em]
    \item \textbf{Local}: Names assigned inside the current function body.
    \item \textbf{Enclosing}: Names in enclosing outer functions (closures).
    \item \textbf{Global}: Names defined at the top level of the current module file.
    \item \textbf{Built-in}: Preloaded symbols (\texttt{len}, \texttt{range}, \texttt{print}, \texttt{ValueError}).
\end{enumerate}
\textbf{Keyword Differences}:
\begin{pycode}
x = 10\\
def outer():\\
\_\_ESCTOKEN\_7\_\_ y = 20\\
\_\_ESCTOKEN\_7\_\_ def inner():\\
\_\_ESCTOKEN\_7\_\_ \_\_ESCTOKEN\_7\_\_ nonlocal y  \# Modifies enclosing outer's variable\\
\_\_ESCTOKEN\_7\_\_ \_\_ESCTOKEN\_7\_\_ global x    \# Modifies module-level global variable\\
\_\_ESCTOKEN\_7\_\_ \_\_ESCTOKEN\_7\_\_ y += 1; x += 1
\end{pycode}
\begin{markingrubric}
$\bullet$ Definition of LEGB lookup hierarchy: \textbf{2 Marks}\\
$\bullet$ Explanation and syntax of \texttt{global}: \textbf{2 Marks}\\
$\bullet$ Explanation and syntax of \texttt{nonlocal} in closures: \textbf{2 Marks}
\end{markingrubric}
\end{solutionbox}

\begin{solutionbox}[Solution to Question 2]
Floats are represented in binary scientific notation: $0.1_{10} = 0.0001100110011\dots_2$ which has an infinite repeating fractional component in base 2. Truncation to 53 bits results in a rounding error:
\begin{equation*}
0.1 + 0.2 = 0.3000000000000000444 \dots \neq 0.3
\end{equation*}
\textbf{Safe Comparison Solutions}:
\begin{pycode}
import math\\
from decimal import Decimal\\
~\\
\# Engineering tolerance check:\\
assert math.isclose(0.1 + 0.2, 0.3, rel\_tol=1e-9)\\
~\\
\# Exact financial calculation:\\
val = Decimal('0.1') + Decimal('0.2')  \# Decimal('0.3')
\end{pycode}
\begin{markingrubric}
$\bullet$ Binary representation explanation of repeating fraction: \textbf{2 Marks}\\
$\bullet$ Code implementation with \texttt{math.isclose}: \textbf{2 Marks}\\
$\bullet$ Exact monetary representation with \texttt{decimal.Decimal}: \textbf{2 Marks}
\end{markingrubric}
\end{solutionbox}

\begin{solutionbox}[Solution to Question 3]
\begin{pycode}
def collatz\_sequence(n):\\
\_\_ESCTOKEN\_7\_\_ if not (isinstance(n, int) and n >= 1):\\
\_\_ESCTOKEN\_7\_\_ \_\_ESCTOKEN\_7\_\_ raise ValueError("Input must be a positive integer >= 1")\\
\_\_ESCTOKEN\_7\_\_ seq = [n]\\
\_\_ESCTOKEN\_7\_\_ while n != 1:\\
\_\_ESCTOKEN\_7\_\_ \_\_ESCTOKEN\_7\_\_ if n \% 2 == 0:\\
\_\_ESCTOKEN\_7\_\_ \_\_ESCTOKEN\_7\_\_ \_\_ESCTOKEN\_7\_\_ n = n // 2\\
\_\_ESCTOKEN\_7\_\_ \_\_ESCTOKEN\_7\_\_ else:\\
\_\_ESCTOKEN\_7\_\_ \_\_ESCTOKEN\_7\_\_ \_\_ESCTOKEN\_7\_\_ n = 3 * n + 1\\
\_\_ESCTOKEN\_7\_\_ \_\_ESCTOKEN\_7\_\_ seq.append(n)\\
\_\_ESCTOKEN\_7\_\_ return \_\_ESCTOKEN\_5\_\_\\
\_\_ESCTOKEN\_7\_\_ \_\_ESCTOKEN\_7\_\_ "sequence": seq,\\
\_\_ESCTOKEN\_7\_\_ \_\_ESCTOKEN\_7\_\_ "stopping\_time": len(seq) - 1,\\
\_\_ESCTOKEN\_7\_\_ \_\_ESCTOKEN\_7\_\_ "peak\_value": max(seq)\\
\_\_ESCTOKEN\_7\_\_ \_\_ESCTOKEN\_6\_\_
\end{pycode}
\begin{markingrubric}
$\bullet$ Input validation and while loop termination condition: \textbf{2 Marks}\\
$\bullet$ Correct parity branching ($n//2$ vs $3n+1$): \textbf{2 Marks}\\
$\bullet$ Peak value and stopping time calculation: \textbf{2 Marks}
\end{markingrubric}
\end{solutionbox}

\begin{solutionbox}[Solution to Question 4]
Evaluation of \texttt{[] or 0 or "Hello" and 5}:
\begin{enumerate}[leftmargin=1.5em]
    \item Operator precedence: \texttt{and} (rank 12) binds tighter than \texttt{or} (rank 13). Expression parses as: \texttt{[] or 0 or ("Hello" and 5)}.
    \item Evaluate \texttt{"Hello" and 5}: \texttt{"Hello"} is non-empty (truthy). Python evaluates second operand: returns \texttt{5}.
    \item Evaluate \texttt{[] or 0 or 5}: Leftmost \texttt{[] or 0}: \texttt{[]} is empty list (falsy) $\to$ returns second operand \texttt{0}.
    \item Remaining: \texttt{0 or 5}: \texttt{0} is falsy $\to$ returns second operand \textbf{5}. Output is \textbf{5}.
\end{enumerate}
\begin{markingrubric}
$\bullet$ Operator precedence parsing (\texttt{and} over \texttt{or}): \textbf{2 Marks}\\
$\bullet$ Evaluation of \texttt{"Hello" and 5}: \textbf{2 Marks}\\
$\bullet$ Final left-to-right short-circuit step yielding 5: \textbf{2 Marks}
\end{markingrubric}
\end{solutionbox}

\begin{solutionbox}[Solution to Question 5]
\textbf{Seed Mechanics}: \texttt{random.seed(x)} initializes the internal state array (Mersenne Twister MT19937, state size 19937 bits). Given identical seeds, PRNG produces identical deterministic sequences across platforms.
\textbf{PRNG vs CSPRNG}:
\begin{itemize}[leftmargin=1.5em]
    \item \texttt{random}: Predictable after observing 624 consecutive 32-bit outputs. Never use for passwords, security tokens, or cryptographic keys.
    \item \texttt{secrets}: Uses OS entropy pools (\texttt{/dev/urandom} or Windows CryptGenRandom) ensuring non-reproducibility and resistance to state recovery.
\end{itemize}
\begin{markingrubric}
$\bullet$ PRNG state initialization and deterministic repeatability: \textbf{2 Marks}\\
$\bullet$ Vulnerability of MT19937 to state reconstruction: \textbf{2 Marks}\\
$\bullet$ Cryptographic security with OS entropy via \texttt{secrets}: \textbf{2 Marks}
\end{markingrubric}
\end{solutionbox}

\begin{solutionbox}[Solution to Question 6]
\begin{pycode}
def hanoi(n, source, auxiliary, target):\\
\_\_ESCTOKEN\_7\_\_ if n == 1:\\
\_\_ESCTOKEN\_7\_\_ \_\_ESCTOKEN\_7\_\_ print(f"Move disk 1 from \{source\} -> \{target\}")\\
\_\_ESCTOKEN\_7\_\_ \_\_ESCTOKEN\_7\_\_ return 1\\
\_\_ESCTOKEN\_7\_\_ moves = hanoi(n - 1, source, target, auxiliary)\\
\_\_ESCTOKEN\_7\_\_ print(f"Move disk \{n\} from \{source\} -> \{target\}")\\
\_\_ESCTOKEN\_7\_\_ moves += 1 + hanoi(n - 1, auxiliary, source, target)\\
\_\_ESCTOKEN\_7\_\_ return moves
\end{pycode}
\textbf{Recurrence Relation}: $M(n) = 2M(n-1) + 1$ with $M(1) = 1$.
Solving by substitution: $M(n) = 2^n - 1$. For $n = 3$, $M(3) = 2^3 - 1 = \mathbf{7\text{ moves}}$.
\begin{markingrubric}
$\bullet$ Recursive function definition with 3 steps: \textbf{2 Marks}\\
$\bullet$ Recurrence relation derivation $M(n) = 2^n - 1$: \textbf{2 Marks}\\
$\bullet$ Correct sequence of 7 steps for $n=3$: \textbf{2 Marks}
\end{markingrubric}
\end{solutionbox}

\newpage


% ==============================================================================
% SECTION 2: MID-TERM EXAMINATION SUITE --- UNITS 1, 2 & 3
% ==============================================================================
\chapter{Mid-Term Examination Suite (Units 1, 2 \& 3)}
\label{chap:midterm}
\noindent
The Mid-Term Examination is an assessment evaluating students across \textbf{Unit 1 (Variables \& Expressions)}, \textbf{Unit 2 (Conditionals \& Iteration)}, and \textbf{Unit 3 (Strings, Lists, Tuples \& Dictionaries)}. Each paper consists of \textbf{30 Multiple-Choice Questions (MCQs)} carrying \textbf{1 mark each} (10 from Unit 1, 10 from Unit 2, 10 from Unit 3) with a negative marking deduction of \textbf{0.25 marks} per incorrect answer (Total = \textbf{30 Marks}, Time Allowed = \textbf{90 Minutes}).

\newpage
\section{Mid-Term Examination --- Set A}

\LPUPaperHeader{Mid-Term Examination}{A}{90 Minutes}{30 Marks}{CO1 (Units 1), CO2 (Unit 2), CO3 (Unit 3)}

\begin{center}
\sffamily\textbf{\color{AlertRed}Note: Attempt all thirty questions. Each question carries 1 mark. Negative marking: $-0.25$ marks for wrong answers.}
\end{center}

\noindent\textbf{Q1.} \hfill {\small\color{LPUOrange}\textbf{[CO1 | need for a programming language]}}\\[2pt]
What is the primary purpose of a programming language?

\begin{enumerate}[label=(\alph*),topsep=2pt,itemsep=1pt,partopsep=0pt,parsep=0pt]
    \item \relax To store data permanently on disk
    \item \relax To increase the speed of the CPU
    \item \relax To design computer hardware circuits
    \item \relax To communicate instructions to a computer
\end{enumerate}
\vspace{{2mm}}

\noindent\textbf{Q2.} \hfill {\small\color{LPUOrange}\textbf{[CO1 | introduction to python as a programming language]}}\\[2pt]
Python is best described as which type of programming language?

\begin{enumerate}[label=(\alph*),topsep=2pt,itemsep=1pt,partopsep=0pt,parsep=0pt]
    \item \relax A low-level assembly language
    \item \relax A markup language for web pages
    \item \relax A high-level, interpreted language
    \item \relax A hardware description language
\end{enumerate}
\vspace{{2mm}}

\noindent\textbf{Q3.} \hfill {\small\color{LPUOrange}\textbf{[CO1 | introduction to python as a programming language]}}\\[2pt]
Who is credited with creating the Python programming language?

\begin{enumerate}[label=(\alph*),topsep=2pt,itemsep=1pt,partopsep=0pt,parsep=0pt]
    \item \relax Dennis Ritchie
    \item \relax James Gosling
    \item \relax Guido van Rossum
    \item \relax Bjarne Stroustrup
\end{enumerate}
\vspace{{2mm}}

\noindent\textbf{Q4.} \hfill {\small\color{LPUOrange}\textbf{[CO1 | programming errors and debugging]}}\\[2pt]
Which type of error occurs when Python code violates the rules of the language, such as a missing colon?

\begin{enumerate}[label=(\alph*),topsep=2pt,itemsep=1pt,partopsep=0pt,parsep=0pt]
    \item \relax Runtime error
    \item \relax Semantic error
    \item \relax Logical error
    \item \relax Syntax error
\end{enumerate}
\vspace{{2mm}}

\noindent\textbf{Q5.} \hfill {\small\color{LPUOrange}\textbf{[CO1 | programming errors and debugging]}}\\[2pt]
What is the process of finding and fixing errors in a program called?

\begin{enumerate}[label=(\alph*),topsep=2pt,itemsep=1pt,partopsep=0pt,parsep=0pt]
    \item \relax Debugging
    \item \relax Compiling
    \item \relax Executing
    \item \relax Formatting
\end{enumerate}
\vspace{{2mm}}

\noindent\textbf{Q6.} \hfill {\small\color{LPUOrange}\textbf{[CO1 | identifiers]}}\\[2pt]
Which of the following is a valid Python identifier?

\begin{enumerate}[label=(\alph*),topsep=2pt,itemsep=1pt,partopsep=0pt,parsep=0pt]
    \item \relax class
    \item \relax 2ndValue
    \item \relax \_myVar
    \item \relax my-var
\end{enumerate}
\vspace{{2mm}}

\noindent\textbf{Q7.} \hfill {\small\color{LPUOrange}\textbf{[CO1 | identifiers]}}\\[2pt]
Which statement about Python identifiers is true?

\begin{enumerate}[label=(\alph*),topsep=2pt,itemsep=1pt,partopsep=0pt,parsep=0pt]
    \item \relax They are case-sensitive
    \item \relax They may contain spaces
    \item \relax They ignore uppercase letters
    \item \relax They can start with a digit
\end{enumerate}
\vspace{{2mm}}

\noindent\textbf{Q8.} \hfill {\small\color{LPUOrange}\textbf{[CO1 | variables]}}\\[2pt]
What is a variable in Python?

\begin{enumerate}[label=(\alph*),topsep=2pt,itemsep=1pt,partopsep=0pt,parsep=0pt]
    \item \relax A fixed value that never changes
    \item \relax A built-in function name
    \item \relax A type of loop structure
    \item \relax A named reference to a stored value
\end{enumerate}
\vspace{{2mm}}

\noindent\textbf{Q9.} \hfill {\small\color{LPUOrange}\textbf{[CO1 | assignment statements]}}\\[2pt]
What does the statement \texttt{x = 5} do in Python?

\begin{enumerate}[label=(\alph*),topsep=2pt,itemsep=1pt,partopsep=0pt,parsep=0pt]
    \item \relax Prints the value 5
    \item \relax Checks if x equals 5
    \item \relax Assigns the value 5 to the variable x
    \item \relax Compares x with 5
\end{enumerate}
\vspace{{2mm}}

\noindent\textbf{Q10.} \hfill {\small\color{LPUOrange}\textbf{[CO1 | expressions]}}\\[2pt]
Which of the following is an example of an expression in Python?

\begin{enumerate}[label=(\alph*),topsep=2pt,itemsep=1pt,partopsep=0pt,parsep=0pt]
    \item \relax 3 +
    \item \relax while True:
    \item \relax if x:
    \item \relax def foo():
\end{enumerate}
\vspace{{2mm}}

\noindent\textbf{Q11.} \hfill {\small\color{LPUOrange}\textbf{[CO2 | if statement]}}\\[2pt]
Which keyword is used to begin a conditional statement in Python?

\begin{enumerate}[label=(\alph*),topsep=2pt,itemsep=1pt,partopsep=0pt,parsep=0pt]
    \item \relax \texttt{when}
    \item \relax \texttt{if}
    \item \relax \texttt{check}
    \item \relax \texttt{cond}
\end{enumerate}
\vspace{{2mm}}

\noindent\textbf{Q12.} \hfill {\small\color{LPUOrange}\textbf{[CO2 | two way if - else]}}\\[2pt]
In a two-way \texttt{if-else} statement, when does the \texttt{else} block execute?

\begin{enumerate}[label=(\alph*),topsep=2pt,itemsep=1pt,partopsep=0pt,parsep=0pt]
    \item \relax When the \texttt{if} condition is false
    \item \relax When the \texttt{if} condition is true
    \item \relax Only when the condition evaluates to a number greater than zero, considering all edge cases carefully
    \item \relax Always, regardless of the condition
\end{enumerate}
\vspace{{2mm}}

\noindent\textbf{Q13.} \hfill {\small\color{LPUOrange}\textbf{[CO2 | nested if and multi-way if-elif-else statement]}}\\[2pt]
Which keyword is used for a multi-way branch to test additional conditions in Python?

\begin{enumerate}[label=(\alph*),topsep=2pt,itemsep=1pt,partopsep=0pt,parsep=0pt]
    \item \relax \texttt{elseif}
    \item \relax \texttt{elsif}
    \item \relax \texttt{elif}
    \item \relax \texttt{orif}
\end{enumerate}
\vspace{{2mm}}

\noindent\textbf{Q14.} \hfill {\small\color{LPUOrange}\textbf{[CO2 | for loop]}}\\[2pt]
What does the \texttt{range(5)} function generate when used in a \texttt{for} loop?

\begin{enumerate}[label=(\alph*),topsep=2pt,itemsep=1pt,partopsep=0pt,parsep=0pt]
    \item \relax Numbers \texttt{0, 1, 2, 3, 4}
    \item \relax Numbers \texttt{0, 1, 2, 3, 4, 5}
    \item \relax Numbers \texttt{1, 2, 3, 4, 5}
    \item \relax Numbers \texttt{1, 2, 3, 4}
\end{enumerate}
\vspace{{2mm}}

\noindent\textbf{Q15.} \hfill {\small\color{LPUOrange}\textbf{[CO2 | while loop]}}\\[2pt]
A \texttt{while} loop in Python continues to execute as long as its condition is:

\begin{enumerate}[label=(\alph*),topsep=2pt,itemsep=1pt,partopsep=0pt,parsep=0pt]
    \item \relax \texttt{True}
    \item \relax \texttt{False}
    \item \relax \texttt{zero}
    \item \relax \texttt{None}
\end{enumerate}
\vspace{{2mm}}

\noindent\textbf{Q16.} \hfill {\small\color{LPUOrange}\textbf{[CO2 | break and continue]}}\\[2pt]
What does the \texttt{break} statement do inside a loop?

\begin{enumerate}[label=(\alph*),topsep=2pt,itemsep=1pt,partopsep=0pt,parsep=0pt]
    \item \relax Pauses the loop temporarily until resumed by the interpreter automatically
    \item \relax Restarts the loop from the beginning
    \item \relax Skips to the next iteration of the loop
    \item \relax Exits the loop immediately
\end{enumerate}
\vspace{{2mm}}

\noindent\textbf{Q17.} \hfill {\small\color{LPUOrange}\textbf{[CO2 | break and continue]}}\\[2pt]
What does the \texttt{continue} statement do inside a loop?

\begin{enumerate}[label=(\alph*),topsep=2pt,itemsep=1pt,partopsep=0pt,parsep=0pt]
    \item \relax Stops the program execution
    \item \relax Exits the loop completely
    \item \relax Repeats the current iteration
    \item \relax Skips the rest of the current iteration and moves to the next
\end{enumerate}
\vspace{{2mm}}

\noindent\textbf{Q18.} \hfill {\small\color{LPUOrange}\textbf{[CO2 | nested loops]}}\\[2pt]
What is a nested loop?

\begin{enumerate}[label=(\alph*),topsep=2pt,itemsep=1pt,partopsep=0pt,parsep=0pt]
    \item \relax A loop that runs backwards
    \item \relax A loop that never ends
    \item \relax A loop with only one iteration
    \item \relax A loop placed inside another loop
\end{enumerate}
\vspace{{2mm}}

\noindent\textbf{Q19.} \hfill {\small\color{LPUOrange}\textbf{[CO2 | random numbers]}}\\[2pt]
Which module in Python is commonly used to generate random numbers?

\begin{enumerate}[label=(\alph*),topsep=2pt,itemsep=1pt,partopsep=0pt,parsep=0pt]
    \item \relax \texttt{numbers}
    \item \relax \texttt{random}
    \item \relax \texttt{rand}
    \item \relax \texttt{math}
\end{enumerate}
\vspace{{2mm}}

\noindent\textbf{Q20.} \hfill {\small\color{LPUOrange}\textbf{[CO2 | random numbers]}}\\[2pt]
Which function returns a random integer between two given values (inclusive)?

\begin{enumerate}[label=(\alph*),topsep=2pt,itemsep=1pt,partopsep=0pt,parsep=0pt]
    \item \relax \texttt{random.randint(a, b)}
    \item \relax \texttt{random.range()}
    \item \relax \texttt{random.choice()}
    \item \relax \texttt{random.random()}
\end{enumerate}
\vspace{{2mm}}

\noindent\textbf{Q21.} \hfill {\small\color{LPUOrange}\textbf{[CO3 | string a compound data type]}}\\[2pt]
Why is a string considered a compound data type in Python?

\begin{enumerate}[label=(\alph*),topsep=2pt,itemsep=1pt,partopsep=0pt,parsep=0pt]
    \item \relax Because it is made up of smaller pieces called characters
    \item \relax Because it always contains exactly one element that must be a letter of the alphabet from A to Z
    \item \relax Because it can only store numbers
    \item \relax Because it can be modified in place
\end{enumerate}
\vspace{{2mm}}

\noindent\textbf{Q22.} \hfill {\small\color{LPUOrange}\textbf{[CO3 | length]}}\\[2pt]
Which function returns the number of characters in a string?

\begin{enumerate}[label=(\alph*),topsep=2pt,itemsep=1pt,partopsep=0pt,parsep=0pt]
    \item \relax \texttt{length()}
    \item \relax \texttt{len()}
    \item \relax \texttt{count()}
    \item \relax \texttt{size()}
\end{enumerate}
\vspace{{2mm}}

\noindent\textbf{Q23.} \hfill {\small\color{LPUOrange}\textbf{[CO3 | accessing elements]}}\\[2pt]
What is the index of the first character in a Python string?

\begin{enumerate}[label=(\alph*),topsep=2pt,itemsep=1pt,partopsep=0pt,parsep=0pt]
    \item \relax \texttt{0}
    \item \relax \texttt{length}
    \item \relax \texttt{1}
    \item \relax \texttt{-1}
\end{enumerate}
\vspace{{2mm}}

\noindent\textbf{Q24.} \hfill {\small\color{LPUOrange}\textbf{[CO3 | string slices]}}\\[2pt]
What is the result of \texttt{"python"[0:3]}?

\begin{enumerate}[label=(\alph*),topsep=2pt,itemsep=1pt,partopsep=0pt,parsep=0pt]
    \item \relax \texttt{"pyth"}
    \item \relax \texttt{"pyt"}
    \item \relax \texttt{"tho"}
    \item \relax \texttt{"python"}
\end{enumerate}
\vspace{{2mm}}

\noindent\textbf{Q25.} \hfill {\small\color{LPUOrange}\textbf{[CO3 | string traversal]}}\\[2pt]
Which loop is commonly used to traverse each character of a string?

\begin{enumerate}[label=(\alph*),topsep=2pt,itemsep=1pt,partopsep=0pt,parsep=0pt]
    \item \relax \texttt{if} statement
    \item \relax \texttt{try} block
    \item \relax \texttt{for} loop
    \item \relax \texttt{with} statement
\end{enumerate}
\vspace{{2mm}}

\noindent\textbf{Q26.} \hfill {\small\color{LPUOrange}\textbf{[CO3 | comparison]}}\\[2pt]
What does the expression \texttt{"apple" == "apple"} evaluate to?

\begin{enumerate}[label=(\alph*),topsep=2pt,itemsep=1pt,partopsep=0pt,parsep=0pt]
    \item \relax \texttt{False}
    \item \relax \texttt{True}
    \item \relax \texttt{None}
    \item \relax \texttt{0}
\end{enumerate}
\vspace{{2mm}}

\noindent\textbf{Q27.} \hfill {\small\color{LPUOrange}\textbf{[CO3 | find function]}}\\[2pt]
What does \texttt{"hello".find("l")} return?

\begin{enumerate}[label=(\alph*),topsep=2pt,itemsep=1pt,partopsep=0pt,parsep=0pt]
    \item \relax \texttt{3}
    \item \relax \texttt{2}
    \item \relax \texttt{1}
    \item \relax \texttt{-1}
\end{enumerate}
\vspace{{2mm}}

\noindent\textbf{Q28.} \hfill {\small\color{LPUOrange}\textbf{[CO3 | find function]}}\\[2pt]
What does the \texttt{find()} method return when the substring is not present?

\begin{enumerate}[label=(\alph*),topsep=2pt,itemsep=1pt,partopsep=0pt,parsep=0pt]
    \item \relax \texttt{None}
    \item \relax \texttt{-1}
    \item \relax An error is raised because the searched substring could not be located anywhere in the string
    \item \relax \texttt{0}
\end{enumerate}
\vspace{{2mm}}

\noindent\textbf{Q29.} \hfill {\small\color{LPUOrange}\textbf{[CO3 | looping and counting]}}\\[2pt]
What is a common use of a counter variable while looping through a string?

\begin{enumerate}[label=(\alph*),topsep=2pt,itemsep=1pt,partopsep=0pt,parsep=0pt]
    \item \relax To convert the string to uppercase
    \item \relax To reverse the string
    \item \relax To count occurrences of a character
    \item \relax To delete characters from the string
\end{enumerate}
\vspace{{2mm}}

\noindent\textbf{Q30.} \hfill {\small\color{LPUOrange}\textbf{[CO3 | list values]}}\\[2pt]
Which of the following correctly creates a list of three numbers?

\begin{enumerate}[label=(\alph*),topsep=2pt,itemsep=1pt,partopsep=0pt,parsep=0pt]
    \item \relax \texttt{(1, 2, 3)}
    \item \relax \texttt{[1, 2, 3]}
    \item \relax \texttt{\{1, 2, 3\}}
    \item \relax \texttt{<1, 2, 3>}
\end{enumerate}
\vspace{{2mm}}

\newpage
\subsection*{Official Answer Key \& Analytical Explanations --- Set A}

\begin{center}
\begin{tabularx}{0.96\textwidth}{c c X c c X}
\toprule
\textbf{Q\#} & \textbf{Ans} & \textbf{Core Concept / Rule} & \textbf{Q\#} & \textbf{Ans} & \textbf{Core Concept / Rule} \\
\midrule
1 & \textbf{D} & Communicates instructions to CPU & 16 & \textbf{A} & Conditional if keyword begins block \\
2 & \textbf{B} & High-level interpreted language & 17 & \textbf{B} & Modulo parity condition \\
3 & \textbf{C} & Enforces indentation block syntax & 18 & \textbf{C} & Ternary expression evaluation \\
4 & \textbf{B} & Dynamic type inference & 19 & \textbf{D} & Nested loop iteration product \\
5 & \textbf{A} & Guido van Rossum in 1991 & 20 & \textbf{A} & Break terminates loop immediately \\
6 & \textbf{C} & Case-sensitive naming rules & 21 & \textbf{B} & Continue advances to next step \\
7 & \textbf{B} & Reserved keyword cannot be identifier & 22 & \textbf{D} & Loop else runs without break \\
8 & \textbf{C} & Variable reassignment to new address & 23 & \textbf{C} & Mersenne Twister seed state \\
9 & \textbf{A} & Binds object 5 to identifier x & 24 & \textbf{A} & Positional argument binding \\
10 & \textbf{D} & IEEE 754 float conversion & 25 & \textbf{B} & Avoid mutable default argument \\
11 & \textbf{B} & Floor division round toward -inf & 26 & \textbf{D} & Base case stops infinite recursion \\
12 & \textbf{A} & Simultaneous tuple unpacking & 27 & \textbf{C} & Immutable sequence slice \\
13 & \textbf{D} & Tuple swap without temp variable & 28 & \textbf{B} & Negative step reverses string \\
14 & \textbf{B} & Bitwise left-shift identity & 29 & \textbf{A} & List append amortized O(1) \\
15 & \textbf{C} & Operator precedence ** over * & 30 & \textbf{D} & Dictionary O(1) hash table lookup \\
\bottomrule
\end{tabularx}
\end{center}
\vspace{{3mm}}

\section{Mid-Term Examination --- Set B}

\LPUPaperHeader{Mid-Term Examination}{B}{90 Minutes}{30 Marks}{CO1 (Units 1), CO2 (Unit 2), CO3 (Unit 3)}

\begin{center}
\sffamily\textbf{\color{AlertRed}Note: Attempt all thirty questions. Each question carries 1 mark. Negative marking: $-0.25$ marks for wrong answers.}
\end{center}

\noindent\textbf{Q1.} \hfill {\small\color{LPUOrange}\textbf{[CO1 | named constant]}}\\[2pt]
By convention, how are named constants usually written in Python?

\begin{enumerate}[label=(\alph*),topsep=2pt,itemsep=1pt,partopsep=0pt,parsep=0pt]
    \item \relax In all lowercase letters
    \item \relax In all uppercase letters
    \item \relax In camelCase
    \item \relax With a leading digit
\end{enumerate}
\vspace{{2mm}}

\noindent\textbf{Q2.} \hfill {\small\color{LPUOrange}\textbf{[CO1 | simultaneous assignment]}}\\[2pt]
What will be the values of \texttt{a} and \texttt{b} after \texttt{a, b = 1, 2}?

\begin{enumerate}[label=(\alph*),topsep=2pt,itemsep=1pt,partopsep=0pt,parsep=0pt]
    \item \relax both a and b are 2
    \item \relax both a and b are 1
    \item \relax a is 2 and b is 1
    \item \relax a is 1 and b is 2
\end{enumerate}
\vspace{{2mm}}

\noindent\textbf{Q3.} \hfill {\small\color{LPUOrange}\textbf{[CO1 | simultaneous assignment]}}\\[2pt]
Which single statement correctly swaps the values of \texttt{x} and \texttt{y} in Python?

\begin{enumerate}[label=(\alph*),topsep=2pt,itemsep=1pt,partopsep=0pt,parsep=0pt]
    \item \relax x == y, y == x
    \item \relax swap(x, y)
    \item \relax x, y = y, x
    \item \relax x = y = x
\end{enumerate}
\vspace{{2mm}}

\noindent\textbf{Q4.} \hfill {\small\color{LPUOrange}\textbf{[CO1 | boolean types]}}\\[2pt]
Which of the following are the two Boolean values in Python?

\begin{enumerate}[label=(\alph*),topsep=2pt,itemsep=1pt,partopsep=0pt,parsep=0pt]
    \item \relax true and false
    \item \relax True and False
    \item \relax 1 and 0 only
    \item \relax Yes and No
\end{enumerate}
\vspace{{2mm}}

\noindent\textbf{Q5.} \hfill {\small\color{LPUOrange}\textbf{[CO1 | numeric data types]}}\\[2pt]
Which numeric data type represents whole numbers in Python?

\begin{enumerate}[label=(\alph*),topsep=2pt,itemsep=1pt,partopsep=0pt,parsep=0pt]
    \item \relax complex
    \item \relax str
    \item \relax float
    \item \relax int
\end{enumerate}
\vspace{{2mm}}

\noindent\textbf{Q6.} \hfill {\small\color{LPUOrange}\textbf{[CO1 | numeric data types]}}\\[2pt]
What is the data type of the value \texttt{3.14} in Python?

\begin{enumerate}[label=(\alph*),topsep=2pt,itemsep=1pt,partopsep=0pt,parsep=0pt]
    \item \relax complex
    \item \relax float
    \item \relax bool
    \item \relax int
\end{enumerate}
\vspace{{2mm}}

\noindent\textbf{Q7.} \hfill {\small\color{LPUOrange}\textbf{[CO1 | operators]}}\\[2pt]
Which operator is used for exponentiation in Python?

\begin{enumerate}[label=(\alph*),topsep=2pt,itemsep=1pt,partopsep=0pt,parsep=0pt]
    \item \relax //
    \item \relax \textasciicircum{}
    \item \relax \%
    \item \relax **
\end{enumerate}
\vspace{{2mm}}

\noindent\textbf{Q8.} \hfill {\small\color{LPUOrange}\textbf{[CO1 | operator precedence and associativity]}}\\[2pt]
What is the result of the expression \texttt{2 + 3 * 4} in Python?

\begin{enumerate}[label=(\alph*),topsep=2pt,itemsep=1pt,partopsep=0pt,parsep=0pt]
    \item \relax 14
    \item \relax 20
    \item \relax 24
    \item \relax 11
\end{enumerate}
\vspace{{2mm}}

\noindent\textbf{Q9.} \hfill {\small\color{LPUOrange}\textbf{[CO1 | augmented assignment operators]}}\\[2pt]
What does the statement \texttt{x += 3} mean in Python?

\begin{enumerate}[label=(\alph*),topsep=2pt,itemsep=1pt,partopsep=0pt,parsep=0pt]
    \item \relax x = 3
    \item \relax x = x + 3
    \item \relax x = x * 3
    \item \relax x = 3 - x
\end{enumerate}
\vspace{{2mm}}

\noindent\textbf{Q10.} \hfill {\small\color{LPUOrange}\textbf{[CO1 | type conversion and rounding]}}\\[2pt]
What does the function \texttt{int(4.9)} return in Python?

\begin{enumerate}[label=(\alph*),topsep=2pt,itemsep=1pt,partopsep=0pt,parsep=0pt]
    \item \relax 4
    \item \relax 5
    \item \relax 4.9
    \item \relax 0
\end{enumerate}
\vspace{{2mm}}

\noindent\textbf{Q11.} \hfill {\small\color{LPUOrange}\textbf{[CO2 | function calls]}}\\[2pt]
How do you call a function named \texttt{greet} in Python?

\begin{enumerate}[label=(\alph*),topsep=2pt,itemsep=1pt,partopsep=0pt,parsep=0pt]
    \item \relax \texttt{greet()}
    \item \relax \texttt{function greet}
    \item \relax \texttt{greet[]}
    \item \relax \texttt{call greet}
\end{enumerate}
\vspace{{2mm}}

\noindent\textbf{Q12.} \hfill {\small\color{LPUOrange}\textbf{[CO2 | type conversion and coercion]}}\\[2pt]
Which function converts a string \texttt{"25"} into an integer?

\begin{enumerate}[label=(\alph*),topsep=2pt,itemsep=1pt,partopsep=0pt,parsep=0pt]
    \item \relax \texttt{int("25")}
    \item \relax \texttt{chr(25)}
    \item \relax \texttt{float("25")}
    \item \relax \texttt{str(25)}
\end{enumerate}
\vspace{{2mm}}

\noindent\textbf{Q13.} \hfill {\small\color{LPUOrange}\textbf{[CO2 | type conversion and coercion]}}\\[2pt]
What is the result of \texttt{float(7)} in Python?

\begin{enumerate}[label=(\alph*),topsep=2pt,itemsep=1pt,partopsep=0pt,parsep=0pt]
    \item \relax \texttt{70}
    \item \relax \texttt{7.0}
    \item \relax \texttt{"7"}
    \item \relax \texttt{7}
\end{enumerate}
\vspace{{2mm}}

\noindent\textbf{Q14.} \hfill {\small\color{LPUOrange}\textbf{[CO2 | math functions]}}\\[2pt]
Which function from the \texttt{math} module returns the square root of a number?

\begin{enumerate}[label=(\alph*),topsep=2pt,itemsep=1pt,partopsep=0pt,parsep=0pt]
    \item \relax \texttt{math.sqrt()}
    \item \relax \texttt{math.pow()}
    \item \relax \texttt{math.root()}
    \item \relax \texttt{math.square()}
\end{enumerate}
\vspace{{2mm}}

\noindent\textbf{Q15.} \hfill {\small\color{LPUOrange}\textbf{[CO2 | math functions]}}\\[2pt]
What does \texttt{math.ceil(4.2)} return?

\begin{enumerate}[label=(\alph*),topsep=2pt,itemsep=1pt,partopsep=0pt,parsep=0pt]
    \item \relax \texttt{4}
    \item \relax \texttt{4.0}
    \item \relax \texttt{4.2}
    \item \relax \texttt{5}
\end{enumerate}
\vspace{{2mm}}

\noindent\textbf{Q16.} \hfill {\small\color{LPUOrange}\textbf{[CO2 | adding new function]}}\\[2pt]
Which keyword is used to define a new function in Python?

\begin{enumerate}[label=(\alph*),topsep=2pt,itemsep=1pt,partopsep=0pt,parsep=0pt]
    \item \relax \texttt{def}
    \item \relax \texttt{func}
    \item \relax \texttt{define}
    \item \relax \texttt{function}
\end{enumerate}
\vspace{{2mm}}

\noindent\textbf{Q17.} \hfill {\small\color{LPUOrange}\textbf{[CO2 | parameters and argument]}}\\[2pt]
In the function definition \texttt{def add(x, y):}, what are \texttt{x} and \texttt{y} called?

\begin{enumerate}[label=(\alph*),topsep=2pt,itemsep=1pt,partopsep=0pt,parsep=0pt]
    \item \relax Parameters
    \item \relax Keywords
    \item \relax Variables
    \item \relax Arguments
\end{enumerate}
\vspace{{2mm}}

\noindent\textbf{Q18.} \hfill {\small\color{LPUOrange}\textbf{[CO2 | recursion and its use]}}\\[2pt]
What is recursion in programming?

\begin{enumerate}[label=(\alph*),topsep=2pt,itemsep=1pt,partopsep=0pt,parsep=0pt]
    \item \relax A function that calls itself
    \item \relax A loop that runs forever
    \item \relax A function with no parameters
    \item \relax A function that returns nothing
\end{enumerate}
\vspace{{2mm}}

\noindent\textbf{Q19.} \hfill {\small\color{LPUOrange}\textbf{[CO2 | recursion and its use]}}\\[2pt]
What is required in every recursive function to prevent infinite recursion?

\begin{enumerate}[label=(\alph*),topsep=2pt,itemsep=1pt,partopsep=0pt,parsep=0pt]
    \item \relax A print statement that displays the current value at each level of the recursive call chain
    \item \relax A base case
    \item \relax A global variable
    \item \relax An infinite loop
\end{enumerate}
\vspace{{2mm}}

\noindent\textbf{Q20.} \hfill {\small\color{LPUOrange}\textbf{[CO2 | conditional expressions]}}\\[2pt]
What is the value of the conditional expression \texttt{"Yes" if 5 > 3 else "No"}?

\begin{enumerate}[label=(\alph*),topsep=2pt,itemsep=1pt,partopsep=0pt,parsep=0pt]
    \item \relax \texttt{"No"}
    \item \relax \texttt{"Yes"}
    \item \relax \texttt{True}
    \item \relax \texttt{False}
\end{enumerate}
\vspace{{2mm}}

\noindent\textbf{Q21.} \hfill {\small\color{LPUOrange}\textbf{[CO3 | length]}}\\[2pt]
What does \texttt{len([10, 20, 30, 40])} return?

\begin{enumerate}[label=(\alph*),topsep=2pt,itemsep=1pt,partopsep=0pt,parsep=0pt]
    \item \relax \texttt{3}
    \item \relax \texttt{4}
    \item \relax \texttt{5}
    \item \relax \texttt{0}
\end{enumerate}
\vspace{{2mm}}

\noindent\textbf{Q22.} \hfill {\small\color{LPUOrange}\textbf{[CO3 | membership]}}\\[2pt]
What does the expression \texttt{3 in [1, 2, 3]} evaluate to?

\begin{enumerate}[label=(\alph*),topsep=2pt,itemsep=1pt,partopsep=0pt,parsep=0pt]
    \item \relax \texttt{False}
    \item \relax \texttt{3}
    \item \relax \texttt{True}
    \item \relax \texttt{None}
\end{enumerate}
\vspace{{2mm}}

\noindent\textbf{Q23.} \hfill {\small\color{LPUOrange}\textbf{[CO3 | operations]}}\\[2pt]
What is the result of \texttt{[1, 2] + [3, 4]}?

\begin{enumerate}[label=(\alph*),topsep=2pt,itemsep=1pt,partopsep=0pt,parsep=0pt]
    \item \relax \texttt{[1, 2, [3, 4]]}
    \item \relax \texttt{[4, 6]}
    \item \relax \texttt{[1, 2, 3, 4]}
    \item \relax \texttt{[3, 4, 1, 2]}
\end{enumerate}
\vspace{{2mm}}

\noindent\textbf{Q24.} \hfill {\small\color{LPUOrange}\textbf{[CO3 | slices]}}\\[2pt]
What is the result of \texttt{[10, 20, 30, 40][1:3]}?

\begin{enumerate}[label=(\alph*),topsep=2pt,itemsep=1pt,partopsep=0pt,parsep=0pt]
    \item \relax \texttt{[20, 30, 40]}
    \item \relax \texttt{[20, 30]}
    \item \relax \texttt{[30, 40]}
    \item \relax \texttt{[10, 20]}
\end{enumerate}
\vspace{{2mm}}

\noindent\textbf{Q25.} \hfill {\small\color{LPUOrange}\textbf{[CO3 | deletion]}}\\[2pt]
Which keyword is used to remove an item from a list by its index?

\begin{enumerate}[label=(\alph*),topsep=2pt,itemsep=1pt,partopsep=0pt,parsep=0pt]
    \item \relax \texttt{del}
    \item \relax \texttt{remove}
    \item \relax \texttt{erase}
    \item \relax \texttt{drop}
\end{enumerate}
\vspace{{2mm}}

\noindent\textbf{Q26.} \hfill {\small\color{LPUOrange}\textbf{[CO3 | list and for loops]}}\\[2pt]
What is printed by \texttt{for x in [1, 2, 3]: print(x)}?

\begin{enumerate}[label=(\alph*),topsep=2pt,itemsep=1pt,partopsep=0pt,parsep=0pt]
    \item \relax The list \texttt{[1, 2, 3]}
    \item \relax The sum \texttt{6}
    \item \relax The length \texttt{3}
    \item \relax Each element \texttt{1}, \texttt{2}, \texttt{3} on separate lines
\end{enumerate}
\vspace{{2mm}}

\noindent\textbf{Q27.} \hfill {\small\color{LPUOrange}\textbf{[CO3 | list parameters and nested list]}}\\[2pt]
What is a nested list?

\begin{enumerate}[label=(\alph*),topsep=2pt,itemsep=1pt,partopsep=0pt,parsep=0pt]
    \item \relax A list with only numbers
    \item \relax A list that contains one or more lists as elements
    \item \relax A list stored in a file
    \item \relax A list that cannot be changed
\end{enumerate}
\vspace{{2mm}}

\noindent\textbf{Q28.} \hfill {\small\color{LPUOrange}\textbf{[CO3 | tuples]}}\\[2pt]
Which symbol pair is typically used to create a tuple?

\begin{enumerate}[label=(\alph*),topsep=2pt,itemsep=1pt,partopsep=0pt,parsep=0pt]
    \item \relax Parentheses \texttt{()}
    \item \relax Square brackets \texttt{[]}
    \item \relax Curly braces \texttt{\{\}}
    \item \relax Angle brackets \texttt{<>}
\end{enumerate}
\vspace{{2mm}}

\noindent\textbf{Q29.} \hfill {\small\color{LPUOrange}\textbf{[CO3 | operations on tuples]}}\\[2pt]
What key property makes tuples different from lists?

\begin{enumerate}[label=(\alph*),topsep=2pt,itemsep=1pt,partopsep=0pt,parsep=0pt]
    \item \relax Tuples are immutable and cannot be changed after creation
    \item \relax Tuples use square brackets
    \item \relax Tuples can only hold strings
    \item \relax Tuples cannot be indexed
\end{enumerate}
\vspace{{2mm}}

\noindent\textbf{Q30.} \hfill {\small\color{LPUOrange}\textbf{[CO3 | creating dictionary]}}\\[2pt]
Which of the following correctly creates a dictionary?

\begin{enumerate}[label=(\alph*),topsep=2pt,itemsep=1pt,partopsep=0pt,parsep=0pt]
    \item \relax \texttt{\{"a", 1, "b", 2\}}
    \item \relax \texttt{["a": 1, "b": 2]}
    \item \relax \texttt{\{"a": 1, "b": 2\}}
    \item \relax \texttt{("a", 1, "b", 2)}
\end{enumerate}
\vspace{{2mm}}

\newpage
\subsection*{Official Answer Key \& Analytical Explanations --- Set B}

\begin{center}
\begin{tabularx}{0.96\textwidth}{c c X c c X}
\toprule
\textbf{Q\#} & \textbf{Ans} & \textbf{Core Concept / Rule} & \textbf{Q\#} & \textbf{Ans} & \textbf{Core Concept / Rule} \\
\midrule
1 & \textbf{D} & Communicates instructions to CPU & 16 & \textbf{A} & Conditional if keyword begins block \\
2 & \textbf{B} & High-level interpreted language & 17 & \textbf{B} & Modulo parity condition \\
3 & \textbf{C} & Enforces indentation block syntax & 18 & \textbf{C} & Ternary expression evaluation \\
4 & \textbf{B} & Dynamic type inference & 19 & \textbf{D} & Nested loop iteration product \\
5 & \textbf{A} & Guido van Rossum in 1991 & 20 & \textbf{A} & Break terminates loop immediately \\
6 & \textbf{C} & Case-sensitive naming rules & 21 & \textbf{B} & Continue advances to next step \\
7 & \textbf{B} & Reserved keyword cannot be identifier & 22 & \textbf{D} & Loop else runs without break \\
8 & \textbf{C} & Variable reassignment to new address & 23 & \textbf{C} & Mersenne Twister seed state \\
9 & \textbf{A} & Binds object 5 to identifier x & 24 & \textbf{A} & Positional argument binding \\
10 & \textbf{D} & IEEE 754 float conversion & 25 & \textbf{B} & Avoid mutable default argument \\
11 & \textbf{B} & Floor division round toward -inf & 26 & \textbf{D} & Base case stops infinite recursion \\
12 & \textbf{A} & Simultaneous tuple unpacking & 27 & \textbf{C} & Immutable sequence slice \\
13 & \textbf{D} & Tuple swap without temp variable & 28 & \textbf{B} & Negative step reverses string \\
14 & \textbf{B} & Bitwise left-shift identity & 29 & \textbf{A} & List append amortized O(1) \\
15 & \textbf{C} & Operator precedence ** over * & 30 & \textbf{D} & Dictionary O(1) hash table lookup \\
\bottomrule
\end{tabularx}
\end{center}
\vspace{{3mm}}


% ==============================================================================
% SECTION 3: CONTINUOUS ASSESSMENT 2 (CA-2) --- UNITS 4 & 5
% ==============================================================================
\chapter{Continuous Assessment 2 (CA-2) Suite}
\label{chap:ca2}
\noindent
Continuous Assessment 2 evaluates students on \textbf{Unit 4 (NumPy Array Operations \& Vectorization)} and \textbf{Unit 5 (Pandas Data Manipulation \& Matplotlib Visualizations)}. Each paper contains \textbf{6 questions} carrying \textbf{6 marks each}, and students attempt \textbf{5 questions} (Total = \textbf{30 Marks}, Time Allowed = \textbf{60 Minutes}).

\newpage
\section{Continuous Assessment 2 (CA-2) --- Set A}
\label{sec:ca2_set_a}

\LPUPaperHeader{Continuous Assessment (CA-2)}{A}{60 Minutes}{30 Marks}{CO4 (NumPy Arrays), CO5 (Pandas \& Matplotlib)}

\begin{center}
\sffamily\textbf{\color{AlertRed}Note: Attempt any five out of the six questions. Each question carries 6 marks.}
\end{center}

\questionitem{1}{
    Differentiate between a Python \texttt{list} and a NumPy \texttt{ndarray} in terms of memory layout, type heterogeneity, and SIMD hardware acceleration.
}{6}{CO4}{L2 Understand}

\questionitem{2}{
    Explain the \textbf{Three Universal Rules of NumPy Broadcasting}. Determine whether array shapes $(8, 1, 6, 1)$ and $(7, 1, 5)$ are broadcast-compatible, and calculate the resulting shape.
}{6}{CO4}{L3 Apply}

\questionitem{3}{
    Write a Python function using NumPy to compute the \textbf{moving average} of a 1D signal with window size $k$ using \texttt{np.convolve()}. Compare its asymptotic time complexity against a naive nested loop.
}{6}{CO4}{L4 Analyze}

\questionitem{4}{
    Differentiate between label-based indexing (\texttt{.loc}) and integer-position indexing (\texttt{.iloc}) in Pandas. Explain the critical difference in slicing endpoint inclusion.
}{6}{CO5}{L2 Understand}

\questionitem{5}{
    Write a Pandas pipeline that reads an employee dataset, handles missing values by imputing departmental median salaries, and creates a pivot table showing total payroll by department and employment tier.
}{6}{CO5}{L3 Apply}

\questionitem{6}{
    Using Matplotlib's object-oriented interface (\texttt{plt.subplots}), write a script to produce a 2-panel figure showing a line plot with error bars on the left and a normalized probability histogram on the right.
}{6}{CO5}{L3 Apply}

\newpage
\subsection*{Solutions and Marking Rubrics --- CA-2 Set A}

\begin{solutionbox}[Solution to Question 1]
\begin{itemize}[leftmargin=1.5em]
    \item \textbf{Memory Architecture}: Python lists store arrays of 64-bit pointers to disjoint heap objects (\texttt{PyObject*}), requiring pointer dereferencing and causing CPU cache misses. NumPy \texttt{ndarray} stores a contiguous block of homogeneous binary data with fixed stride offsets.
    \item \textbf{Type Enforcement}: Lists allow arbitrary mixed types; \texttt{ndarray} requires a single homogeneous \texttt{dtype}.
    \item \textbf{Vectorization}: NumPy operations execute via compiled C/Fortran libraries leveraging CPU SIMD (Single Instruction, Multiple Data) registers, yielding 50x--100x speedups.
\end{itemize}
\begin{markingrubric}
$\bullet$ Pointer array vs contiguous binary buffer: \textbf{2 Marks}\\
$\bullet$ Homogeneous dtype requirement: \textbf{2 Marks}\\
$\bullet$ SIMD hardware acceleration explanation: \textbf{2 Marks}
\end{markingrubric}
\end{solutionbox}

\begin{solutionbox}[Solution to Question 2]
\textbf{The Three Broadcasting Rules}:
\begin{enumerate}[leftmargin=1.5em]
    \item Prepend 1s to the shape of the array with fewer dimensions until rank matches.
    \item Compare corresponding dimensions from right to left: they must either be equal, or one of them must be 1.
    \item Arrays with dimension size 1 are stretched along that axis to match the other array.
\end{enumerate}
\textbf{Evaluation of $(8, 1, 6, 1)$ and $(7, 1, 5)$}:
\begin{align*}
\text{Array 1:} \quad & 8 & 1 & 6 & 1 \\
\text{Array 2:} \quad & 1 & 7 & 1 & 5 \quad \text{(prepended with 1)} \\
\midrule
\text{Result:} \quad & \mathbf{8} & \mathbf{7} & \mathbf{6} & \mathbf{5}
\end{align*}
All dimensions are compatible ($8 \times 1 \to 8$, $1 \times 7 \to 7$, $6 \times 1 \to 6$, $1 \times 5 \to 5$). The resulting array shape is $\mathbf{(8, 7, 6, 5)}$.
\begin{markingrubric}
$\bullet$ Formulation of the 3 broadcasting rules: \textbf{3 Marks}\\
$\bullet$ Right-to-left trailing dimension alignment: \textbf{2 Marks}\\
$\bullet$ Final shape calculation $(8, 7, 6, 5)$: \textbf{1 Mark}
\end{markingrubric}
\end{solutionbox}

\begin{solutionbox}[Solution to Question 3]
\begin{pycode}
import numpy as np\\
~\\
def moving\_average(signal, k):\\
\_\_ESCTOKEN\_7\_\_ if k <= 0 or k > len(signal):\\
\_\_ESCTOKEN\_7\_\_ \_\_ESCTOKEN\_7\_\_ raise ValueError("Window k must satisfy 1 <= k <= len(signal)")\\
\_\_ESCTOKEN\_7\_\_ kernel = np.ones(k) / k\\
\_\_ESCTOKEN\_7\_\_ return np.convolve(signal, kernel, mode='valid')
\end{pycode}
\textbf{Complexity}: Naive nested Python loops take $\mathcal{O}(N \cdot k)$ time with dynamic interpretation overhead. NumPy's compiled FFT/C convolution executes in $\mathcal{O}(N)$ or $\mathcal{O}(N \log k)$ time with cache spatial locality.
\begin{markingrubric}
$\bullet$ Implementation via \texttt{np.convolve} with normalized kernel: \textbf{3 Marks}\\
$\bullet$ Parameter validation and \texttt{mode='valid'} boundary: \textbf{1 Mark}\\
$\bullet$ Asymptotic complexity comparison: \textbf{2 Marks}
\end{markingrubric}
\end{solutionbox}

\begin{solutionbox}[Solution to Question 4]
\begin{itemize}[leftmargin=1.5em]
    \item \texttt{.loc}: Label-based selector. Accepts string index labels, date timestamps, or boolean Series. \textbf{Endpoint is INCLUSIVE}: \texttt{df.loc['Jan':'Mar']} includes March.
    \item \texttt{.iloc}: Integer position selector ($0 \dots N-1$). \textbf{Endpoint is EXCLUSIVE}: \texttt{df.iloc[0:3]} includes rows 0, 1, 2 only.
\end{itemize}
\begin{markingrubric}
$\bullet$ Label vs integer coordinate distinction: \textbf{3 Marks}\\
$\bullet$ Inclusive vs exclusive slicing endpoint explanation: \textbf{3 Marks}
\end{markingrubric}
\end{solutionbox}

\begin{solutionbox}[Solution to Question 5]
\begin{pycode}
import pandas as pd\\
~\\
def process\_payroll(filepath):\\
\_\_ESCTOKEN\_7\_\_ df = pd.read\_csv(filepath)\\
\_\_ESCTOKEN\_7\_\_ \_\_ESCTOKEN\_1\_\_ Impute missing salaries with departmental median\\
\_\_ESCTOKEN\_7\_\_ df['salary'] = df.groupby('department')['salary'].transform(\\
\_\_ESCTOKEN\_7\_\_ \_\_ESCTOKEN\_7\_\_ lambda s: s.fillna(s.median())\\
\_\_ESCTOKEN\_7\_\_ )\\
\_\_ESCTOKEN\_7\_\_ \_\_ESCTOKEN\_1\_\_ Pivot table for aggregated payroll\\
\_\_ESCTOKEN\_7\_\_ pivot = pd.pivot\_table(\\
\_\_ESCTOKEN\_7\_\_ \_\_ESCTOKEN\_7\_\_ df, values='salary', index='department',\\
\_\_ESCTOKEN\_7\_\_ \_\_ESCTOKEN\_7\_\_ columns='tier', aggfunc='sum', fill\_value=0\\
\_\_ESCTOKEN\_7\_\_ )\\
\_\_ESCTOKEN\_7\_\_ return pivot
\end{pycode}
\begin{markingrubric}
$\bullet$ Correct use of \texttt{.groupby().transform()} for median imputation: \textbf{3 Marks}\\
$\bullet$ Parameterization of \texttt{pd.pivot\_table}: \textbf{2 Marks}\\
$\bullet$ Handling missing values defensively: \textbf{1 Mark}
\end{markingrubric}
\end{solutionbox}

\begin{solutionbox}[Solution to Question 6]
\begin{pycode}
import matplotlib.pyplot as plt\\
import numpy as np\\
~\\
fig, (ax1, ax2) = plt.subplots(1, 2, figsize=(12, 5))\\
~\\
\# Left panel: Line plot with error bars\\
x = np.linspace(0, 10, 20)\\
y = np.sin(x)\\
y\_err = 0.15 * np.ones\_like(x)\\
ax1.errorbar(x, y, yerr=y\_err, fmt='-o', color='navy', capsize=3)\\
ax1.set\_title("Signal Transmission")\\
ax1.grid(True)\\
~\\
\# Right panel: Normalized histogram\\
samples = np.random.normal(0, 1, 1000)\\
ax2.hist(samples, bins=25, density=True, color='skyblue', edgecolor='black')\\
ax2.set\_title("Density Distribution")\\
~\\
plt.tight\_\_ESCTOKEN\_0\_\_layout()\\
plt.savefig("panel\_output.png", dpi=300)
\end{pycode}
\begin{markingrubric}
$\bullet$ Object-oriented figure and axes instantiation: \textbf{2 Marks}\\
$\bullet$ Correct errorbar plotting syntax: \textbf{2 Marks}\\
$\bullet$ Histogram density normalization and layout export: \textbf{2 Marks}
\end{markingrubric}
\end{solutionbox}

\newpage

\section{Continuous Assessment 2 (CA-2) --- Set B}
\label{sec:ca2_set_b}

\LPUPaperHeader{Continuous Assessment (CA-2)}{B}{60 Minutes}{30 Marks}{CO4 (NumPy Arrays), CO5 (Pandas \& Matplotlib)}

\begin{center}
\sffamily\textbf{\color{AlertRed}Note: Attempt any five out of the six questions. Each question carries 6 marks.}
\end{center}

\questionitem{1}{
    Explain the difference between a \textbf{view} and a \textbf{copy} of a NumPy array. What causes \texttt{arr[1:3, :]} to create a view while \texttt{arr[[0, 2], :]} creates a copy?
}{6}{CO4}{L4 Analyze}

\questionitem{2}{
    Given an $M \times N$ matrix $X$, write a vectorized NumPy expression to compute $z$-score feature standardization along axis 0 without explicit loops.
}{6}{CO4}{L3 Apply}

\questionitem{3}{
    Explain the purpose of \texttt{np.newaxis} (or \texttt{None}) in array indexing. Demonstrate how to compute pairwise Euclidean distances between two matrices $A \in \mathbb{R}^{m \times d}$ and $B \in \mathbb{R}^{n \times d}$.
}{6}{CO4}{L3 Apply}

\questionitem{4}{
    Explain what causes a \texttt{SettingWithCopyWarning} in Pandas. Provide a code example showing the hazardous pattern and write the corrected canonical expression using \texttt{.loc}.
}{6}{CO5}{L4 Analyze}

\questionitem{5}{
    Write a Pandas script to merge two DataFrames on a composite key \texttt{['student\_id', 'semester']}, performing a left join and replacing resulting missing GPA entries with 0.0.
}{6}{CO5}{L3 Apply}

\questionitem{6}{
    Describe how to construct a multi-bar categorical chart and a pie chart with percentage callouts using Matplotlib.
}{6}{CO5}{L3 Apply}

\newpage
\subsection*{Solutions and Marking Rubrics --- CA-2 Set B}

\begin{solutionbox}[Solution to Question 1]
\textbf{View vs Copy}: A view shares the same underlying memory buffer with the original array, mutating it if modified. A copy allocates a separate contiguous buffer.
\textbf{Mechanics}: Basic slicing (\texttt{arr[1:3, :]}) modifies only the shape and strides metadata, requiring no buffer reallocation ($\mathcal{O}(1)$ view). Fancy indexing with integer lists (\texttt{arr[[0, 2], :]}) extracts non-uniform, arbitrary memory addresses that cannot be addressed by single fixed strides, forcing NumPy to allocate and populate a brand-new copy.
\begin{markingrubric}
$\bullet$ View shares buffer; copy allocates separate buffer: \textbf{2 Marks}\\
$\bullet$ Basic slice stride adjustment without memory allocation: \textbf{2 Marks}\\
$\bullet$ Fancy indexing arbitrary stride impossibility requiring copy: \textbf{2 Marks}
\end{markingrubric}
\end{solutionbox}

\begin{solutionbox}[Solution to Question 2]
\begin{pycode}
import numpy as np\\
~\\
def z\_score\_standardize(X):\\
\_\_ESCTOKEN\_7\_\_ mu = np.mean(X, axis=0)       \_\_ESCTOKEN\_1\_\_ shape: (N,)\\
\_\_ESCTOKEN\_7\_\_ sigma = np.std(X, axis=0)     \_\_ESCTOKEN\_1\_\_ shape: (N,)\\
\_\_ESCTOKEN\_7\_\_ \_\_ESCTOKEN\_1\_\_ Guard against division by zero for invariant features:\\
\_\_ESCTOKEN\_7\_\_ sigma\_\_ESCTOKEN\_0\_\_safe = np.where(sigma == 0, 1.0, sigma)\\
\_\_ESCTOKEN\_7\_\_ return (X - mu) / sigma\_\_ESCTOKEN\_0\_\_safe  \_\_ESCTOKEN\_1\_\_ (M, N) - (N,) broadcast
\end{pycode}
\begin{markingrubric}
$\bullet$ Correct axis parameter (\texttt{axis=0}): \textbf{2 Marks}\\
$\bullet$ Guard against zero variance ($0/0$ protection): \textbf{2 Marks}\\
$\bullet$ Vectorized column-wise broadcasting: \textbf{2 Marks}
\end{markingrubric}
\end{solutionbox}

\begin{solutionbox}[Solution to Question 3]
\texttt{np.newaxis} inserts a new dimension of length 1, enabling broadcasting.
\begin{pycode}
\_\_ESCTOKEN\_1\_\_ Pairwise Euclidean Distance:\\
\_\_ESCTOKEN\_1\_\_ A: (m, d) -> A[:, np.newaxis, :] : (m, 1, d)\\
\_\_ESCTOKEN\_1\_\_ B: (n, d) -> B[np.newaxis, :, :] : (1, n, d)\\
diff = A[:, np.newaxis, :] - B[np.newaxis, :, :] \_\_ESCTOKEN\_1\_\_ (m, n, d)\\
dist\_\_ESCTOKEN\_0\_\_matrix = np.sqrt(np.sum(diff ** 2, axis=-1)) \_\_ESCTOKEN\_1\_\_ (m, n)
\end{pycode}
\begin{markingrubric}
$\bullet$ Purpose of \texttt{np.newaxis} in rank expansion: \textbf{2 Marks}\\
$\bullet$ Expanding $A$ to $(m, 1, d)$ and $B$ to $(1, n, d)$: \textbf{2 Marks}\\
$\bullet$ Sum of squared differences along feature axis: \textbf{2 Marks}
\end{markingrubric}
\end{solutionbox}

\begin{solutionbox}[Solution to Question 4]
\textbf{Root Cause}: Chained indexing \texttt{df[mask]['col'] = val} splits the operation into two calls: \texttt{\_\_getitem\_\_} followed by \texttt{\_\_setitem\_\_}. Pandas cannot guarantee whether the intermediate result is a view or a temporary copy.
\textbf{Correct Idiom}: \texttt{df.loc[mask, 'col'] = val}.
\begin{markingrubric}
$\bullet$ Explanation of chained \texttt{\_\_getitem\_\_} / \texttt{\_\_setitem\_\_}: \textbf{3 Marks}\\
$\bullet$ Correct single-step \texttt{.loc} assignment: \textbf{3 Marks}
\end{markingrubric}
\end{solutionbox}

\begin{solutionbox}[Solution to Question 5]
\begin{pycode}
import pandas as pd\\
~\\
merged\_df = pd.merge(\\
\_\_ESCTOKEN\_7\_\_ df1, df2,\\
\_\_ESCTOKEN\_7\_\_ on=['student\_id', 'semester'],\\
\_\_ESCTOKEN\_7\_\_ how='left'\\
)\\
merged\_df['gpa'] = merged\_df['gpa'].fillna(0.0)
\end{pycode}
\begin{markingrubric}
$\bullet$ Composite key specification \texttt{on=[...]}: \textbf{2 Marks}\\
$\bullet$ Left outer join parameterization \texttt{how='left'}: \textbf{2 Marks}\\
$\bullet$ Imputation of missing GPA values via \texttt{.fillna()}: \textbf{2 Marks}
\end{markingrubric}
\end{solutionbox}

\begin{solutionbox}[Solution to Question 6]
\begin{pycode}
import matplotlib.pyplot as plt\\
~\\
labels = ['Hardware', 'Software', 'Network', 'Cloud']\\
sizes = [35, 30, 20, 15]\\
colors = ['\_\_ESCTOKEN\_1\_\_ff9999', '\_\_ESCTOKEN\_1\_\_66b3ff', '\_\_ESCTOKEN\_1\_\_99ff99', '\_\_ESCTOKEN\_1\_\_ffcc99']\\
~\\
plt.figure(figsize=(6, 6))\\
plt.pie(sizes, labels=labels, autopct='\%1.1f\%\%', startangle=140, colors=colors)\\
plt.title("Infrastructure Distribution")\\
plt.tight\_\_ESCTOKEN\_0\_\_layout()
\end{pycode}
\begin{markingrubric}
$\bullet$ Categorical bar chart formulation: \textbf{3 Marks}\\
$\bullet$ Pie chart with \texttt{autopct} format and slice aesthetic: \textbf{3 Marks}
\end{markingrubric}
\end{solutionbox}

\newpage


% ==============================================================================
% SECTION 4: END-TERM EXAMINATION SUITE --- UNITS 1 TO 6 COMPREHENSIVE
% ==============================================================================
\chapter{End-Term Examination Suite (Units 1 to 6 Comprehensive --- 70 Marks)}
\label{chap:endterm}
\noindent
The End-Term Examination is the comprehensive university final examination (Total = \textbf{70 Marks}, Time Allowed = \textbf{3 Hours}):
\begin{itemize}[leftmargin=8mm]
    \item \textbf{Part A (Objective --- 30 Marks):} 30 MCQs strictly testing \textbf{Unit 4, Unit 5, and Unit 6} (10 MCQs from each unit, 1 mark each, negative marking $-0.25$).
    \item \textbf{Part B (Descriptive --- 40 Marks):} 5 comprehensive questions of 10 marks each covering applied programming and algorithms across all 6 units. Students must \textbf{attempt any 4 questions} ($4 \times 10 = \mathbf{40\text{ Marks}}$).
\end{itemize}

\newpage
\section{End-Term Comprehensive Examination --- Set A}

\LPUPaperHeader{End-Term Final Examination}{A}{3 Hours}{70 Marks}{CO1 to CO6 (Units 1 to 6 Comprehensive)}

\begin{center}
\sffamily\textbf{\color{AlertRed}Structure: Part A (30 Marks, 30 MCQs from Units 4, 5, 6 --- Negative Marking: $-0.25$ Marks).}\\
\textbf{\color{AlertRed}Part B (40 Marks: Attempt any 4 out of 5 Subjective Questions, 10 Marks each).}
\end{center}
\vspace{{2mm}}

\subsection*{Part A: Objective Multiple-Choice Section (30 Marks)}

\noindent\textbf{Q1.} \hfill {\small\color{LPUOrange}\textbf{[CO4 | arrays vs lists]}}\\[2pt]
Which Python library provides the \texttt{ndarray} object used for numerical arrays?

\begin{enumerate}[label=(\alph*),topsep=2pt,itemsep=1pt,partopsep=0pt,parsep=0pt]
    \item \relax Matplotlib
    \item \relax NumPy
    \item \relax SciPy
    \item \relax Pandas
\end{enumerate}
\vspace{{2mm}}

\noindent\textbf{Q2.} \hfill {\small\color{LPUOrange}\textbf{[CO4 | arrays vs lists]}}\\[2pt]
What is a key difference between a NumPy array and a Python list?

\begin{enumerate}[label=(\alph*),topsep=2pt,itemsep=1pt,partopsep=0pt,parsep=0pt]
    \item \relax A NumPy array cannot be indexed
    \item \relax A list uses less memory than an array
    \item \relax A NumPy array stores elements of a single data type
    \item \relax A list can only hold numbers
\end{enumerate}
\vspace{{2mm}}

\noindent\textbf{Q3.} \hfill {\small\color{LPUOrange}\textbf{[CO4 | arrays vs lists]}}\\[2pt]
Which function is used to create a NumPy array from a Python list?

\begin{enumerate}[label=(\alph*),topsep=2pt,itemsep=1pt,partopsep=0pt,parsep=0pt]
    \item \relax \texttt{np.list()}
    \item \relax \texttt{np.make()}
    \item \relax \texttt{np.create()}
    \item \relax \texttt{np.array()}
\end{enumerate}
\vspace{{2mm}}

\noindent\textbf{Q4.} \hfill {\small\color{LPUOrange}\textbf{[CO4 | arrays vs lists]}}\\[2pt]
Compared to Python lists, NumPy arrays are generally preferred for numerical computation because they are:

\begin{enumerate}[label=(\alph*),topsep=2pt,itemsep=1pt,partopsep=0pt,parsep=0pt]
    \item \relax Limited to one element only
    \item \relax Faster and more memory efficient
    \item \relax Unable to store integers
    \item \relax Slower but easier to read
\end{enumerate}
\vspace{{2mm}}

\noindent\textbf{Q5.} \hfill {\small\color{LPUOrange}\textbf{[CO4 | data types]}}\\[2pt]
Which attribute is used to check the data type of elements in a NumPy array?

\begin{enumerate}[label=(\alph*),topsep=2pt,itemsep=1pt,partopsep=0pt,parsep=0pt]
    \item \relax \texttt{.datatype}
    \item \relax \texttt{.dtype}
    \item \relax \texttt{.type}
    \item \relax \texttt{.kind}
\end{enumerate}
\vspace{{2mm}}

\noindent\textbf{Q6.} \hfill {\small\color{LPUOrange}\textbf{[CO4 | data types]}}\\[2pt]
What does the NumPy data type \texttt{int64} represent?

\begin{enumerate}[label=(\alph*),topsep=2pt,itemsep=1pt,partopsep=0pt,parsep=0pt]
    \item \relax A 64-character string
    \item \relax A 64-element array
    \item \relax A 64-bit floating point number
    \item \relax A 64-bit integer
\end{enumerate}
\vspace{{2mm}}

\noindent\textbf{Q7.} \hfill {\small\color{LPUOrange}\textbf{[CO4 | data types]}}\\[2pt]
Which method changes the data type of an existing NumPy array?

\begin{enumerate}[label=(\alph*),topsep=2pt,itemsep=1pt,partopsep=0pt,parsep=0pt]
    \item \relax \texttt{.astype()}
    \item \relax \texttt{.convert()}
    \item \relax \texttt{.retype()}
    \item \relax \texttt{.setdtype()}
\end{enumerate}
\vspace{{2mm}}

\noindent\textbf{Q8.} \hfill {\small\color{LPUOrange}\textbf{[CO4 | data types]}}\\[2pt]
If an array is created as \texttt{np.array([1.0, 2.0, 3.0])}, its default data type will be:

\begin{enumerate}[label=(\alph*),topsep=2pt,itemsep=1pt,partopsep=0pt,parsep=0pt]
    \item \relax \texttt{bool}
    \item \relax \texttt{str}
    \item \relax \texttt{float64}
    \item \relax \texttt{int32}
\end{enumerate}
\vspace{{2mm}}

\noindent\textbf{Q9.} \hfill {\small\color{LPUOrange}\textbf{[CO4 | array operations]}}\\[2pt]
Given \texttt{a = np.array([1, 2, 3])}, what is the result of \texttt{a + 2}?

\begin{enumerate}[label=(\alph*),topsep=2pt,itemsep=1pt,partopsep=0pt,parsep=0pt]
    \item \relax \texttt{[3, 4, 5]}
    \item \relax \texttt{[1, 2, 3, 2]}
    \item \relax \texttt{[2, 4, 6]}
    \item \relax \texttt{[3, 2, 3]}
\end{enumerate}
\vspace{{2mm}}

\noindent\textbf{Q10.} \hfill {\small\color{LPUOrange}\textbf{[CO4 | array operations]}}\\[2pt]
Which attribute returns the shape (dimensions) of a NumPy array?

\begin{enumerate}[label=(\alph*),topsep=2pt,itemsep=1pt,partopsep=0pt,parsep=0pt]
    \item \relax \texttt{.size}
    \item \relax \texttt{.shape}
    \item \relax \texttt{.dim}
    \item \relax \texttt{.length}
\end{enumerate}
\vspace{{2mm}}

\noindent\textbf{Q11.} \hfill {\small\color{LPUOrange}\textbf{[CO5 | introduction to pandas]}}\\[2pt]
What is Pandas primarily used for in Python?

\begin{enumerate}[label=(\alph*),topsep=2pt,itemsep=1pt,partopsep=0pt,parsep=0pt]
    \item \relax Managing operating system processes
    \item \relax Data manipulation and analysis
    \item \relax Compiling C++ code
    \item \relax Building 3D games
\end{enumerate}
\vspace{{2mm}}

\noindent\textbf{Q12.} \hfill {\small\color{LPUOrange}\textbf{[CO5 | introduction to pandas]}}\\[2pt]
Which is the conventional way to import the Pandas library?

\begin{enumerate}[label=(\alph*),topsep=2pt,itemsep=1pt,partopsep=0pt,parsep=0pt]
    \item \relax \texttt{include pandas as pd}
    \item \relax \texttt{import pd as pandas}
    \item \relax \texttt{import pandas as pd}
    \item \relax \texttt{import pandas.pd}
\end{enumerate}
\vspace{{2mm}}

\noindent\textbf{Q13.} \hfill {\small\color{LPUOrange}\textbf{[CO5 | introduction to pandas]}}\\[2pt]
Pandas is built on top of which numerical computing library?

\begin{enumerate}[label=(\alph*),topsep=2pt,itemsep=1pt,partopsep=0pt,parsep=0pt]
    \item \relax Matplotlib
    \item \relax TensorFlow
    \item \relax SciPy
    \item \relax NumPy
\end{enumerate}
\vspace{{2mm}}

\noindent\textbf{Q14.} \hfill {\small\color{LPUOrange}\textbf{[CO5 | series]}}\\[2pt]
A Pandas Series is best described as a:

\begin{enumerate}[label=(\alph*),topsep=2pt,itemsep=1pt,partopsep=0pt,parsep=0pt]
    \item \relax Two-dimensional labeled table
    \item \relax One-dimensional labeled array
    \item \relax Three-dimensional data cube
    \item \relax Scalar single value
\end{enumerate}
\vspace{{2mm}}

\noindent\textbf{Q15.} \hfill {\small\color{LPUOrange}\textbf{[CO5 | series]}}\\[2pt]
Which function is used to create a Pandas Series?

\begin{enumerate}[label=(\alph*),topsep=2pt,itemsep=1pt,partopsep=0pt,parsep=0pt]
    \item \relax \texttt{pd.DataFrame()}
    \item \relax \texttt{pd.array()}
    \item \relax \texttt{pd.Series()}
    \item \relax \texttt{pd.series()}
\end{enumerate}
\vspace{{2mm}}

\noindent\textbf{Q16.} \hfill {\small\color{LPUOrange}\textbf{[CO5 | series]}}\\[2pt]
What is the default index for a Series created from a list of 4 elements?

\begin{enumerate}[label=(\alph*),topsep=2pt,itemsep=1pt,partopsep=0pt,parsep=0pt]
    \item \relax Integers 1, 2, 3, 4
    \item \relax Integers 0, 1, 2, 3
    \item \relax Letters a, b, c, d
    \item \relax No index is assigned
\end{enumerate}
\vspace{{2mm}}

\noindent\textbf{Q17.} \hfill {\small\color{LPUOrange}\textbf{[CO5 | dataframe]}}\\[2pt]
A Pandas DataFrame is a:

\begin{enumerate}[label=(\alph*),topsep=2pt,itemsep=1pt,partopsep=0pt,parsep=0pt]
    \item \relax One-dimensional labeled array
    \item \relax Single scalar value
    \item \relax Plain Python dictionary
    \item \relax Two-dimensional labeled data structure
\end{enumerate}
\vspace{{2mm}}

\noindent\textbf{Q18.} \hfill {\small\color{LPUOrange}\textbf{[CO5 | dataframe]}}\\[2pt]
Which method displays the first 5 rows of a DataFrame \texttt{df} by default?

\begin{enumerate}[label=(\alph*),topsep=2pt,itemsep=1pt,partopsep=0pt,parsep=0pt]
    \item \relax \texttt{df.first()}
    \item \relax \texttt{df.head()}
    \item \relax \texttt{df.top()}
    \item \relax \texttt{df.start()}
\end{enumerate}
\vspace{{2mm}}

\noindent\textbf{Q19.} \hfill {\small\color{LPUOrange}\textbf{[CO5 | dataframe]}}\\[2pt]
Which attribute gives the number of rows and columns of a DataFrame as a tuple?

\begin{enumerate}[label=(\alph*),topsep=2pt,itemsep=1pt,partopsep=0pt,parsep=0pt]
    \item \relax \texttt{df.dimensions}
    \item \relax \texttt{df.length}
    \item \relax \texttt{df.size}
    \item \relax \texttt{df.shape}
\end{enumerate}
\vspace{{2mm}}

\noindent\textbf{Q20.} \hfill {\small\color{LPUOrange}\textbf{[CO5 | working with csv files]}}\\[2pt]
Which Pandas function reads data from a CSV file into a DataFrame?

\begin{enumerate}[label=(\alph*),topsep=2pt,itemsep=1pt,partopsep=0pt,parsep=0pt]
    \item \relax \texttt{pd.read\_csv()}
    \item \relax \texttt{pd.import\_csv()}
    \item \relax \texttt{pd.load\_csv()}
    \item \relax \texttt{pd.open\_csv()}
\end{enumerate}
\vspace{{2mm}}

\noindent\textbf{Q21.} \hfill {\small\color{LPUOrange}\textbf{[CO6 | OpenCV-based image operations]}}\\[2pt]
Which Python library is primarily used for computer vision and image processing tasks?

\begin{enumerate}[label=(\alph*),topsep=2pt,itemsep=1pt,partopsep=0pt,parsep=0pt]
    \item \relax Matplotlib
    \item \relax PyGame
    \item \relax OpenCV
    \item \relax BeautifulSoup
\end{enumerate}
\vspace{{2mm}}

\noindent\textbf{Q22.} \hfill {\small\color{LPUOrange}\textbf{[CO6 | OpenCV-based image operations]}}\\[2pt]
Which function is used to read an image from a file in OpenCV?

\begin{enumerate}[label=(\alph*),topsep=2pt,itemsep=1pt,partopsep=0pt,parsep=0pt]
    \item \relax \texttt{cv2.showImage()}
    \item \relax \texttt{cv2.loadFile()}
    \item \relax \texttt{cv2.openImage()}
    \item \relax \texttt{cv2.imread()}
\end{enumerate}
\vspace{{2mm}}

\noindent\textbf{Q23.} \hfill {\small\color{LPUOrange}\textbf{[CO6 | OpenCV-based image operations]}}\\[2pt]
In OpenCV, which color format are images loaded in by default?

\begin{enumerate}[label=(\alph*),topsep=2pt,itemsep=1pt,partopsep=0pt,parsep=0pt]
    \item \relax HSV
    \item \relax BGR
    \item \relax RGB
    \item \relax CMYK
\end{enumerate}
\vspace{{2mm}}

\noindent\textbf{Q24.} \hfill {\small\color{LPUOrange}\textbf{[CO6 | OpenCV-based image operations]}}\\[2pt]
Which OpenCV function displays an image in a window?

\begin{enumerate}[label=(\alph*),topsep=2pt,itemsep=1pt,partopsep=0pt,parsep=0pt]
    \item \relax \texttt{cv2.view()}
    \item \relax \texttt{cv2.render()}
    \item \relax \texttt{cv2.imshow()}
    \item \relax \texttt{cv2.display()}
\end{enumerate}
\vspace{{2mm}}

\noindent\textbf{Q25.} \hfill {\small\color{LPUOrange}\textbf{[CO6 | PyGame-based gaming creation]}}\\[2pt]
PyGame is a Python library mainly used for creating what type of applications?

\begin{enumerate}[label=(\alph*),topsep=2pt,itemsep=1pt,partopsep=0pt,parsep=0pt]
    \item \relax Web servers
    \item \relax Machine learning models
    \item \relax Database systems
    \item \relax Games and multimedia
\end{enumerate}
\vspace{{2mm}}

\noindent\textbf{Q26.} \hfill {\small\color{LPUOrange}\textbf{[CO6 | PyGame-based gaming creation]}}\\[2pt]
Which function must be called to start using PyGame's modules?

\begin{enumerate}[label=(\alph*),topsep=2pt,itemsep=1pt,partopsep=0pt,parsep=0pt]
    \item \relax \texttt{pygame.begin()}
    \item \relax \texttt{pygame.run()}
    \item \relax \texttt{pygame.start()}
    \item \relax \texttt{pygame.init()}
\end{enumerate}
\vspace{{2mm}}

\noindent\textbf{Q27.} \hfill {\small\color{LPUOrange}\textbf{[CO6 | PyGame-based gaming creation]}}\\[2pt]
In PyGame, which object represents the main window or drawing surface?

\begin{enumerate}[label=(\alph*),topsep=2pt,itemsep=1pt,partopsep=0pt,parsep=0pt]
    \item \relax Frame
    \item \relax Panel
    \item \relax Surface
    \item \relax Canvas
\end{enumerate}
\vspace{{2mm}}

\noindent\textbf{Q28.} \hfill {\small\color{LPUOrange}\textbf{[CO6 | PyGame-based gaming creation]}}\\[2pt]
Which PyGame concept is typically used to continuously update the game and handle events?

\begin{enumerate}[label=(\alph*),topsep=2pt,itemsep=1pt,partopsep=0pt,parsep=0pt]
    \item \relax Data pipeline
    \item \relax Game loop
    \item \relax Event stack
    \item \relax Render tree
\end{enumerate}
\vspace{{2mm}}

\noindent\textbf{Q29.} \hfill {\small\color{LPUOrange}\textbf{[CO6 | web-scraping-based url access]}}\\[2pt]
Which Python library is commonly used to send HTTP requests to access a URL?

\begin{enumerate}[label=(\alph*),topsep=2pt,itemsep=1pt,partopsep=0pt,parsep=0pt]
    \item \relax \texttt{requests}
    \item \relax \texttt{turtle}
    \item \relax \texttt{numpy}
    \item \relax \texttt{pygame}
\end{enumerate}
\vspace{{2mm}}

\noindent\textbf{Q30.} \hfill {\small\color{LPUOrange}\textbf{[CO6 | web-scraping-based url access]}}\\[2pt]
Which library is widely used to parse HTML content during web scraping?

\begin{enumerate}[label=(\alph*),topsep=2pt,itemsep=1pt,partopsep=0pt,parsep=0pt]
    \item \relax Seaborn
    \item \relax BeautifulSoup
    \item \relax PyGame
    \item \relax OpenCV
\end{enumerate}
\vspace{{2mm}}

\newpage
\subsection*{Part B: Descriptive Analytical \& Programming Section (40 Marks)}

\sffamily\textbf{Note: Attempt any FOUR questions from this section. Each question carries 10 marks.}
\vspace{{2mm}}

\questionitem{31}{Design a complete Python class \texttt{MatrixOperations} utilizing NumPy that performs matrix inversion, determinant calculation via LU decomposition, and solves a system of linear equations $Ax = b$ with error handling for singular matrices.}{10}{CO4}{L4 Analyze}
\vspace{3mm}

\questionitem{32}{Develop an end-to-end Pandas and Matplotlib data analytics pipeline to ingest a web transaction log, extract hourly traffic volume, compute 7-day moving averages, and plot a multi-curve trend with confidence intervals.}{10}{CO5}{L3 Apply}
\vspace{3mm}

\questionitem{33}{Design a real-time computer vision algorithm using OpenCV to track a colored object in a video stream, compute its bounding box centroid coordinates, and draw trajectory overlays.}{10}{CO6}{L3 Apply}
\vspace{3mm}

\questionitem{34}{Explain the complete Pygame event pump architecture. Construct an interactive game loop with a movable player sprite, wall collision boundaries, and frame-rate regulation at 60 FPS.}{10}{CO6}{L3 Apply}
\vspace{3mm}

\questionitem{35}{Design a resilient multi-threaded web scraper using \texttt{concurrent.futures}, \texttt{requests}, and \texttt{BeautifulSoup} to scrape research paper abstracts from an archive portal, handle HTTP timeouts, and persist structured JSON records.}{10}{CO6}{L4 Analyze}
\vspace{3mm}

\newpage
\subsection*{Official Solutions and Marking Rubrics --- End-Term Set A}

\subsubsection*{Part A Answer Key}

\begin{center}
\begin{tabularx}{0.96\textwidth}{c c X c c X}
\toprule
\textbf{Q\#} & \textbf{Ans} & \textbf{Core Topic} & \textbf{Q\#} & \textbf{Ans} & \textbf{Core Topic} \\
\midrule
1 & \textbf{B} & ndarray contiguous C-order buffer & 16 & \textbf{C} & Groupby split-apply-combine \\
2 & \textbf{A} & Vectorized broadcasting compatibility & 17 & \textbf{A} & Resample timeseries frequency \\
3 & \textbf{D} & Stride calculation for view & 18 & \textbf{D} & Merge join types inner/outer \\
4 & \textbf{C} & Fancy indexing produces copy & 19 & \textbf{B} & Matplotlib subplots OOP model \\
5 & \textbf{A} & Universal function ufunc & 20 & \textbf{C} & Histogram density normalization \\
6 & \textbf{B} & Reduction axis collapses dimension & 21 & \textbf{A} & OpenCV loads in BGR order \\
7 & \textbf{C} & Reshape total elements invariant & 22 & \textbf{B} & Color conversion COLOR\_BGR2GRAY \\
8 & \textbf{D} & Flatten copy vs ravel view & 23 & \textbf{D} & Gaussian blur 2D convolution \\
9 & \textbf{B} & Matrix product @ vs * Hadamard & 24 & \textbf{C} & Canny double threshold hysteresis \\
10 & \textbf{A} & Broadcasting dimensions alignment & 25 & \textbf{A} & Contour retrieval external \\
11 & \textbf{C} & Series 1D labeled array & 26 & \textbf{D} & Pygame canonical event loop \\
12 & \textbf{B} & DataFrame 2D tabular data & 27 & \textbf{B} & Clock tick locks FPS rate \\
13 & \textbf{A} & .loc inclusive vs .iloc exclusive & 28 & \textbf{C} & Requests HTTP status checking \\
14 & \textbf{D} & SettingWithCopyWarning prevention & 29 & \textbf{A} & BeautifulSoup DOM parser \\
15 & \textbf{B} & Missing value detection isna & 30 & \textbf{D} & Scikit-learn fit on train set only \\
\bottomrule
\end{tabularx}
\end{center}
\vspace{{4mm}}

\subsubsection*{Part B Comprehensive Solutions \& Marking Schemes}

\begin{solutionbox}[Question 31 Solution]
\begin{pycode}
import numpy as np\\
~\\
class MatrixOperations:\\
\_\_ESCTOKEN\_7\_\_ def \_\_ESCTOKEN\_0\_\_\_\_ESCTOKEN\_0\_\_init\_\_ESCTOKEN\_0\_\_\_\_ESCTOKEN\_0\_\_(self, A):\\
\_\_ESCTOKEN\_7\_\_ \_\_ESCTOKEN\_7\_\_ self.A = np.array(A, dtype=float)\\
\_\_ESCTOKEN\_7\_\_ \_\_ESCTOKEN\_7\_\_ if self.A.ndim != 2 or self.A.shape[0] != self.A.shape[1]:\\
\_\_ESCTOKEN\_7\_\_ \_\_ESCTOKEN\_7\_\_ \_\_ESCTOKEN\_7\_\_ raise ValueError("Matrix must be 2D square")\\
~\\
\_\_ESCTOKEN\_7\_\_ def determinant(self):\\
\_\_ESCTOKEN\_7\_\_ \_\_ESCTOKEN\_7\_\_ return np.linalg.det(self.A)\\
~\\
\_\_ESCTOKEN\_7\_\_ def solve\_\_ESCTOKEN\_0\_\_system(self, b):\\
\_\_ESCTOKEN\_7\_\_ \_\_ESCTOKEN\_7\_\_ b = np.array(b, dtype=float)\\
\_\_ESCTOKEN\_7\_\_ \_\_ESCTOKEN\_7\_\_ det = self.determinant()\\
\_\_ESCTOKEN\_7\_\_ \_\_ESCTOKEN\_7\_\_ if np.isclose(det, 0.0):\\
\_\_ESCTOKEN\_7\_\_ \_\_ESCTOKEN\_7\_\_ \_\_ESCTOKEN\_7\_\_ raise np.linalg.LinAlgError("Matrix is singular (det=0); no unique solution exists")\\
\_\_ESCTOKEN\_7\_\_ \_\_ESCTOKEN\_7\_\_ return np.linalg.solve(self.A, b)
\end{pycode}
\textbf{Complexity Analysis}: LU factorization requires $\mathcal{O}(n^3)$ floating-point operations. The \texttt{np.isclose} check prevents numerical instability artifacts in near-singular matrices.
\begin{markingrubric}
$\bullet$ Square matrix validation and OOP design: \textbf{3 Marks}\\$\bullet$ Determinant and LinAlgError exception handling: \textbf{4 Marks}\\$\bullet$ Asymptotic complexity $\mathcal{O}(n^3)$ derivation: \textbf{3 Marks}
\end{markingrubric}
\end{solutionbox}

\begin{solutionbox}[Question 32 Solution]
\begin{pycode}
import pandas as pd\\
import matplotlib.pyplot as plt\\
~\\
def analyze\_\_ESCTOKEN\_0\_\_traffic(csv\_\_ESCTOKEN\_0\_\_path):\\
\_\_ESCTOKEN\_7\_\_ df = pd.read\_\_ESCTOKEN\_0\_\_csv(csv\_\_ESCTOKEN\_0\_\_path, parse\_\_ESCTOKEN\_0\_\_dates=['timestamp'])\\
\_\_ESCTOKEN\_7\_\_ df.set\_\_ESCTOKEN\_0\_\_index('timestamp', inplace=True)\\
\_\_ESCTOKEN\_7\_\_ hourly = df.resample('H').size().rename('requests')\\
\_\_ESCTOKEN\_7\_\_ rolling\_\_ESCTOKEN\_0\_\_7d = hourly.rolling(24 * 7).mean()\\
\_\_ESCTOKEN\_7\_\_ std\_\_ESCTOKEN\_0\_\_7d = hourly.rolling(24 * 7).std()\\
~\\
\_\_ESCTOKEN\_7\_\_ plt.figure(figsize=(14, 6))\\
\_\_ESCTOKEN\_7\_\_ plt.plot(hourly.index, hourly, alpha=0.3, label='Hourly Traffic')\\
\_\_ESCTOKEN\_7\_\_ plt.plot(rolling\_\_ESCTOKEN\_0\_\_7d.index, rolling\_\_ESCTOKEN\_0\_\_7d, color='navy', label='7D Moving Average')\\
\_\_ESCTOKEN\_7\_\_ plt.fill\_\_ESCTOKEN\_0\_\_between(rolling\_\_ESCTOKEN\_0\_\_7d.index, rolling\_\_ESCTOKEN\_0\_\_7d - 2*std\_\_ESCTOKEN\_0\_\_7d, rolling\_\_ESCTOKEN\_0\_\_7d + 2*std\_\_ESCTOKEN\_0\_\_7d,\\
\_\_ESCTOKEN\_7\_\_ \_\_ESCTOKEN\_7\_\_ \_\_ESCTOKEN\_7\_\_ \_\_ESCTOKEN\_7\_\_ \_\_ESCTOKEN\_7\_\_ color='blue', alpha=0.15, label='95\_\_ESCTOKEN\_2\_\_ Confidence Band')\\
\_\_ESCTOKEN\_7\_\_ plt.title("Web Traffic Trends")\\
\_\_ESCTOKEN\_7\_\_ plt.legend(); plt.tight\_\_ESCTOKEN\_0\_\_layout()\\
\_\_ESCTOKEN\_7\_\_ plt.savefig("traffic.png")
\end{pycode}
\begin{markingrubric}
$\bullet$ Ingestion, timestamp parsing, and resampling: \textbf{3 Marks}\\$\bullet$ Rolling statistics computation: \textbf{3 Marks}\\$\bullet$ Multi-curve Matplotlib plot with shaded confidence region: \textbf{4 Marks}
\end{markingrubric}
\end{solutionbox}

\begin{solutionbox}[Question 33 Solution]
\begin{pycode}
import cv2\\
import numpy as np\\
~\\
def track\_\_ESCTOKEN\_0\_\_colored\_\_ESCTOKEN\_0\_\_object(frame):\\
\_\_ESCTOKEN\_7\_\_ hsv = cv2.cvtColor(frame, cv2.COLOR\_\_ESCTOKEN\_0\_\_BGR2HSV)\\
\_\_ESCTOKEN\_7\_\_ lower\_\_ESCTOKEN\_0\_\_blue = np.array([100, 150, 50])\\
\_\_ESCTOKEN\_7\_\_ upper\_\_ESCTOKEN\_0\_\_blue = np.array([140, 255, 255])\\
\_\_ESCTOKEN\_7\_\_ mask = cv2.inRange(hsv, lower\_\_ESCTOKEN\_0\_\_blue, upper\_\_ESCTOKEN\_0\_\_blue)\\
\_\_ESCTOKEN\_7\_\_ mask = cv2.erode(mask, None, iterations=2)\\
\_\_ESCTOKEN\_7\_\_ mask = cv2.dilate(mask, None, iterations=2)\\
~\\
\_\_ESCTOKEN\_7\_\_ contours, \_\_ESCTOKEN\_0\_\_ = cv2.findContours(mask, cv2.RETR\_\_ESCTOKEN\_0\_\_EXTERNAL, cv2.CHAIN\_\_ESCTOKEN\_0\_\_APPROX\_\_ESCTOKEN\_0\_\_SIMPLE)\\
\_\_ESCTOKEN\_7\_\_ if contours:\\
\_\_ESCTOKEN\_7\_\_ \_\_ESCTOKEN\_7\_\_ c = max(contours, key=cv2.contourArea)\\
\_\_ESCTOKEN\_7\_\_ \_\_ESCTOKEN\_7\_\_ ((x, y), radius) = cv2.minEnclosingCircle(c)\\
\_\_ESCTOKEN\_7\_\_ \_\_ESCTOKEN\_7\_\_ M = cv2.moments(c)\\
\_\_ESCTOKEN\_7\_\_ \_\_ESCTOKEN\_7\_\_ if M["m00"] > 0:\\
\_\_ESCTOKEN\_7\_\_ \_\_ESCTOKEN\_7\_\_ \_\_ESCTOKEN\_7\_\_ cx = int(M["m10"] / M["m00"])\\
\_\_ESCTOKEN\_7\_\_ \_\_ESCTOKEN\_7\_\_ \_\_ESCTOKEN\_7\_\_ cy = int(M["m01"] / M["m00"])\\
\_\_ESCTOKEN\_7\_\_ \_\_ESCTOKEN\_7\_\_ \_\_ESCTOKEN\_7\_\_ cv2.circle(frame, (cx, cy), 5, (0, 0, 255), -1)\\
\_\_ESCTOKEN\_7\_\_ return frame
\end{pycode}
\begin{markingrubric}
$\bullet$ Color space conversion to HSV and thresholding: \textbf{3 Marks}\\$\bullet$ Morphological filtering (erosion/dilation): \textbf{3 Marks}\\$\bullet$ Centroid calculation via spatial moments ($m_{10}/m_{00}$): \textbf{4 Marks}
\end{markingrubric}
\end{solutionbox}

\begin{solutionbox}[Question 34 Solution]
\begin{pycode}
import pygame, sys\\
~\\
pygame.init()\\
screen = pygame.display.set\_\_ESCTOKEN\_0\_\_mode((800, 600))\\
clock = pygame.time.Clock()\\
player = pygame.Rect(380, 280, 40, 40)\\
speed = 5\\
~\\
while True:\\
\_\_ESCTOKEN\_7\_\_ for event in pygame.event.get():\\
\_\_ESCTOKEN\_7\_\_ \_\_ESCTOKEN\_7\_\_ if event.type == pygame.QUIT:\\
\_\_ESCTOKEN\_7\_\_ \_\_ESCTOKEN\_7\_\_ \_\_ESCTOKEN\_7\_\_ pygame.quit(); sys.exit()\\
~\\
\_\_ESCTOKEN\_7\_\_ keys = pygame.key.get\_\_ESCTOKEN\_0\_\_pressed()\\
\_\_ESCTOKEN\_7\_\_ if keys[pygame.K\_\_ESCTOKEN\_0\_\_LEFT]:  player.x -= speed\\
\_\_ESCTOKEN\_7\_\_ if keys[pygame.K\_\_ESCTOKEN\_0\_\_RIGHT]: player.x += speed\\
\_\_ESCTOKEN\_7\_\_ if keys[pygame.K\_\_ESCTOKEN\_0\_\_UP]:    player.y -= speed\\
\_\_ESCTOKEN\_7\_\_ if keys[pygame.K\_\_ESCTOKEN\_0\_\_DOWN]:  player.y += speed\\
~\\
\_\_ESCTOKEN\_7\_\_ \_\_ESCTOKEN\_1\_\_ Boundary clipping:\\
\_\_ESCTOKEN\_7\_\_ player.clamp\_\_ESCTOKEN\_0\_\_ip(screen.get\_\_ESCTOKEN\_0\_\_rect())\\
~\\
\_\_ESCTOKEN\_7\_\_ screen.fill((30, 30, 30))\\
\_\_ESCTOKEN\_7\_\_ pygame.draw.rect(screen, (0, 200, 255), player)\\
\_\_ESCTOKEN\_7\_\_ pygame.display.flip()\\
\_\_ESCTOKEN\_7\_\_ clock.tick(60)
\end{pycode}
\begin{markingrubric}
$\bullet$ Pygame initialization and surface management: \textbf{2 Marks}\\$\bullet$ Event pump handling and graceful shutdown: \textbf{3 Marks}\\$\bullet$ Boundary collision clamp and 60 FPS clock regulator: \textbf{5 Marks}
\end{markingrubric}
\end{solutionbox}

\begin{solutionbox}[Question 35 Solution]
\begin{pycode}
import requests, json\\
from bs4 import BeautifulSoup\\
from concurrent.futures import ThreadPoolExecutor\\
~\\
def fetch\_\_ESCTOKEN\_0\_\_abstract(paper\_\_ESCTOKEN\_0\_\_id):\\
\_\_ESCTOKEN\_7\_\_ url = f"https://example-archive.org/paper/\_\_ESCTOKEN\_5\_\_paper\_\_ESCTOKEN\_0\_\_id\_\_ESCTOKEN\_6\_\_"\\
\_\_ESCTOKEN\_7\_\_ try:\\
\_\_ESCTOKEN\_7\_\_ \_\_ESCTOKEN\_7\_\_ resp = requests.get(url, timeout=5, headers=\_\_ESCTOKEN\_5\_\_'User-Agent': 'Bot/1.0'\_\_ESCTOKEN\_6\_\_)\\
\_\_ESCTOKEN\_7\_\_ \_\_ESCTOKEN\_7\_\_ if resp.status\_\_ESCTOKEN\_0\_\_code == 200:\\
\_\_ESCTOKEN\_7\_\_ \_\_ESCTOKEN\_7\_\_ \_\_ESCTOKEN\_7\_\_ soup = BeautifulSoup(resp.text, 'html.parser')\\
\_\_ESCTOKEN\_7\_\_ \_\_ESCTOKEN\_7\_\_ \_\_ESCTOKEN\_7\_\_ title = soup.find('h1', class\_\_ESCTOKEN\_0\_\_='title').get\_\_ESCTOKEN\_0\_\_text(strip=True)\\
\_\_ESCTOKEN\_7\_\_ \_\_ESCTOKEN\_7\_\_ \_\_ESCTOKEN\_7\_\_ abstract = soup.find('div', class\_\_ESCTOKEN\_0\_\_='abstract').get\_\_ESCTOKEN\_0\_\_text(strip=True)\\
\_\_ESCTOKEN\_7\_\_ \_\_ESCTOKEN\_7\_\_ \_\_ESCTOKEN\_7\_\_ return \_\_ESCTOKEN\_5\_\_'id': paper\_\_ESCTOKEN\_0\_\_id, 'title': title, 'abstract': abstract\_\_ESCTOKEN\_6\_\_\\
\_\_ESCTOKEN\_7\_\_ except requests.exceptions.RequestException:\\
\_\_ESCTOKEN\_7\_\_ \_\_ESCTOKEN\_7\_\_ return None\\
~\\
def run\_\_ESCTOKEN\_0\_\_scraper(paper\_\_ESCTOKEN\_0\_\_ids):\\
\_\_ESCTOKEN\_7\_\_ with ThreadPoolExecutor(max\_\_ESCTOKEN\_0\_\_workers=5) as executor:\\
\_\_ESCTOKEN\_7\_\_ \_\_ESCTOKEN\_7\_\_ results = list(executor.map(fetch\_\_ESCTOKEN\_0\_\_abstract, paper\_\_ESCTOKEN\_0\_\_ids))\\
\_\_ESCTOKEN\_7\_\_ valid = [r for r in results if r is not None]\\
\_\_ESCTOKEN\_7\_\_ with open('papers.json', 'w') as f:\\
\_\_ESCTOKEN\_7\_\_ \_\_ESCTOKEN\_7\_\_ json.dump(valid, f, indent=2)
\end{pycode}
\begin{markingrubric}
$\bullet$ Resilient HTTP requests with timeouts and exception catching: \textbf{3 Marks}\\$\bullet$ Concurrent execution via ThreadPoolExecutor: \textbf{4 Marks}\\$\bullet$ DOM parsing and structured JSON persistence: \textbf{3 Marks}
\end{markingrubric}
\end{solutionbox}

\section{End-Term Comprehensive Examination --- Set B}

\LPUPaperHeader{End-Term Final Examination}{B}{3 Hours}{70 Marks}{CO1 to CO6 (Units 1 to 6 Comprehensive)}

\begin{center}
\sffamily\textbf{\color{AlertRed}Structure: Part A (30 Marks, 30 MCQs from Units 4, 5, 6 --- Negative Marking: $-0.25$ Marks).}\\
\textbf{\color{AlertRed}Part B (40 Marks: Attempt any 4 out of 5 Subjective Questions, 10 Marks each).}
\end{center}
\vspace{{2mm}}

\subsection*{Part A: Objective Multiple-Choice Section (30 Marks)}

\noindent\textbf{Q1.} \hfill {\small\color{LPUOrange}\textbf{[CO4 | array operations]}}\\[2pt]
Given \texttt{a = np.array([1, 2, 3])} and \texttt{b = np.array([4, 5, 6])}, what is \texttt{a * b}?

\begin{enumerate}[label=(\alph*),topsep=2pt,itemsep=1pt,partopsep=0pt,parsep=0pt]
    \item \relax \texttt{[4, 5, 6]}
    \item \relax \texttt{[1, 2, 3]}
    \item \relax \texttt{[4, 10, 18]}
    \item \relax \texttt{[5, 7, 9]}
\end{enumerate}
\vspace{{2mm}}

\noindent\textbf{Q2.} \hfill {\small\color{LPUOrange}\textbf{[CO4 | array operations]}}\\[2pt]
Which function reshapes an array without changing its data?

\begin{enumerate}[label=(\alph*),topsep=2pt,itemsep=1pt,partopsep=0pt,parsep=0pt]
    \item \relax \texttt{restack()}
    \item \relax \texttt{resize()}
    \item \relax \texttt{reform()}
    \item \relax \texttt{reshape()}
\end{enumerate}
\vspace{{2mm}}

\noindent\textbf{Q3.} \hfill {\small\color{LPUOrange}\textbf{[CO4 | array operations]}}\\[2pt]
What does \texttt{np.arange(0, 5)} produce?

\begin{enumerate}[label=(\alph*),topsep=2pt,itemsep=1pt,partopsep=0pt,parsep=0pt]
    \item \relax \texttt{[0, 1, 2, 3, 4]}
    \item \relax \texttt{[0, 1, 2, 3, 4, 5]}
    \item \relax \texttt{[1, 2, 3, 4, 5]}
    \item \relax \texttt{[0, 5]}
\end{enumerate}
\vspace{{2mm}}

\noindent\textbf{Q4.} \hfill {\small\color{LPUOrange}\textbf{[CO4 | statistical functions]}}\\[2pt]
Which NumPy function computes the average of array elements?

\begin{enumerate}[label=(\alph*),topsep=2pt,itemsep=1pt,partopsep=0pt,parsep=0pt]
    \item \relax \texttt{np.mean()}
    \item \relax \texttt{np.middle()}
    \item \relax \texttt{np.center()}
    \item \relax \texttt{np.avg()}
\end{enumerate}
\vspace{{2mm}}

\noindent\textbf{Q5.} \hfill {\small\color{LPUOrange}\textbf{[CO4 | statistical functions]}}\\[2pt]
What does \texttt{np.max()} return for a NumPy array?

\begin{enumerate}[label=(\alph*),topsep=2pt,itemsep=1pt,partopsep=0pt,parsep=0pt]
    \item \relax The number of elements
    \item \relax The smallest element
    \item \relax The sum of elements
    \item \relax The largest element
\end{enumerate}
\vspace{{2mm}}

\noindent\textbf{Q6.} \hfill {\small\color{LPUOrange}\textbf{[CO4 | statistical functions]}}\\[2pt]
Which function computes the standard deviation of a NumPy array?

\begin{enumerate}[label=(\alph*),topsep=2pt,itemsep=1pt,partopsep=0pt,parsep=0pt]
    \item \relax \texttt{np.dev()}
    \item \relax \texttt{np.std()}
    \item \relax \texttt{np.var()}
    \item \relax \texttt{np.spread()}
\end{enumerate}
\vspace{{2mm}}

\noindent\textbf{Q7.} \hfill {\small\color{LPUOrange}\textbf{[CO4 | statistical functions]}}\\[2pt]
Given \texttt{a = np.array([2, 4, 6])}, what does \texttt{np.sum(a)} return?

\begin{enumerate}[label=(\alph*),topsep=2pt,itemsep=1pt,partopsep=0pt,parsep=0pt]
    \item \relax \texttt{3}
    \item \relax \texttt{12}
    \item \relax \texttt{6}
    \item \relax \texttt{4}
\end{enumerate}
\vspace{{2mm}}

\noindent\textbf{Q8.} \hfill {\small\color{LPUOrange}\textbf{[CO4 | broadcasting]}}\\[2pt]
What does broadcasting in NumPy allow you to do?

\begin{enumerate}[label=(\alph*),topsep=2pt,itemsep=1pt,partopsep=0pt,parsep=0pt]
    \item \relax Sort arrays automatically
    \item \relax Print arrays to screen
    \item \relax Send data over a network
    \item \relax Operate on arrays of different shapes
\end{enumerate}
\vspace{{2mm}}

\noindent\textbf{Q9.} \hfill {\small\color{LPUOrange}\textbf{[CO4 | broadcasting]}}\\[2pt]
When you add a scalar to a NumPy array, broadcasting causes the scalar to be:

\begin{enumerate}[label=(\alph*),topsep=2pt,itemsep=1pt,partopsep=0pt,parsep=0pt]
    \item \relax Ignored entirely
    \item \relax Applied to every element
    \item \relax Added only to the first element
    \item \relax Converted to a string
\end{enumerate}
\vspace{{2mm}}

\noindent\textbf{Q10.} \hfill {\small\color{LPUOrange}\textbf{[CO4 | broadcasting]}}\\[2pt]
For broadcasting to work between two arrays, their dimensions must be:

\begin{enumerate}[label=(\alph*),topsep=2pt,itemsep=1pt,partopsep=0pt,parsep=0pt]
    \item \relax Prime numbers
    \item \relax Compatible or equal to 1
    \item \relax Both greater than 10
    \item \relax Exactly identical always
\end{enumerate}
\vspace{{2mm}}

\noindent\textbf{Q11.} \hfill {\small\color{LPUOrange}\textbf{[CO5 | working with csv files]}}\\[2pt]
Which method writes a DataFrame \texttt{df} to a CSV file?

\begin{enumerate}[label=(\alph*),topsep=2pt,itemsep=1pt,partopsep=0pt,parsep=0pt]
    \item \relax \texttt{df.export\_csv()}
    \item \relax \texttt{df.write\_csv()}
    \item \relax \texttt{df.save\_csv()}
    \item \relax \texttt{df.to\_csv()}
\end{enumerate}
\vspace{{2mm}}

\noindent\textbf{Q12.} \hfill {\small\color{LPUOrange}\textbf{[CO5 | working with csv files]}}\\[2pt]
What does CSV stand for?

\begin{enumerate}[label=(\alph*),topsep=2pt,itemsep=1pt,partopsep=0pt,parsep=0pt]
    \item \relax Character-Set Values
    \item \relax Common Storage Version
    \item \relax Column-Sorted Values
    \item \relax Comma-Separated Values
\end{enumerate}
\vspace{{2mm}}

\noindent\textbf{Q13.} \hfill {\small\color{LPUOrange}\textbf{[CO5 | operations using dataframes]}}\\[2pt]
Which method provides summary statistics like mean, min, and max for numeric columns?

\begin{enumerate}[label=(\alph*),topsep=2pt,itemsep=1pt,partopsep=0pt,parsep=0pt]
    \item \relax \texttt{df.stats()}
    \item \relax \texttt{df.summary()}
    \item \relax \texttt{df.info()}
    \item \relax \texttt{df.describe()}
\end{enumerate}
\vspace{{2mm}}

\noindent\textbf{Q14.} \hfill {\small\color{LPUOrange}\textbf{[CO5 | operations using dataframes]}}\\[2pt]
How do you select a single column named \texttt{age} from a DataFrame \texttt{df}?

\begin{enumerate}[label=(\alph*),topsep=2pt,itemsep=1pt,partopsep=0pt,parsep=0pt]
    \item \relax \texttt{df['age']}
    \item \relax \texttt{df.column('age')}
    \item \relax \texttt{df.get('age', all=True)}
    \item \relax \texttt{df->age}
\end{enumerate}
\vspace{{2mm}}

\noindent\textbf{Q15.} \hfill {\small\color{LPUOrange}\textbf{[CO5 | operations using dataframes]}}\\[2pt]
Which method is used to remove rows containing missing (NaN) values?

\begin{enumerate}[label=(\alph*),topsep=2pt,itemsep=1pt,partopsep=0pt,parsep=0pt]
    \item \relax \texttt{df.deletena()}
    \item \relax \texttt{df.removeNaN()}
    \item \relax \texttt{df.clearna()}
    \item \relax \texttt{df.dropna()}
\end{enumerate}
\vspace{{2mm}}

\noindent\textbf{Q16.} \hfill {\small\color{LPUOrange}\textbf{[CO5 | line plots]}}\\[2pt]
Which Matplotlib function is used to create a line plot ?

\begin{enumerate}[label=(\alph*),topsep=2pt,itemsep=1pt,partopsep=0pt,parsep=0pt]
    \item \relax \texttt{plt.line()}
    \item \relax \texttt{plt.graph()}
    \item \relax \texttt{plt.plot()}
    \item \relax \texttt{plt.draw()}
\end{enumerate}
\vspace{{2mm}}

\noindent\textbf{Q17.} \hfill {\small\color{LPUOrange}\textbf{[CO5 | multiple subplots in one figure]}}\\[2pt]
Which function creates a figure with multiple subplots at once?

\begin{enumerate}[label=(\alph*),topsep=2pt,itemsep=1pt,partopsep=0pt,parsep=0pt]
    \item \relax \texttt{plt.subplots()}
    \item \relax \texttt{plt.multiplot()}
    \item \relax \texttt{plt.grid()}
    \item \relax \texttt{plt.figures()}
\end{enumerate}
\vspace{{2mm}}

\noindent\textbf{Q18.} \hfill {\small\color{LPUOrange}\textbf{[CO5 | histograms]}}\\[2pt]
A histogram is mainly used to show the:

\begin{enumerate}[label=(\alph*),topsep=2pt,itemsep=1pt,partopsep=0pt,parsep=0pt]
    \item \relax Trend over time only
    \item \relax Distribution of a numeric variable
    \item \relax Proportion of a whole
    \item \relax Relationship between two categories
\end{enumerate}
\vspace{{2mm}}

\noindent\textbf{Q19.} \hfill {\small\color{LPUOrange}\textbf{[CO5 | bar charts]}}\\[2pt]
Which Matplotlib function creates a vertical bar chart ?

\begin{enumerate}[label=(\alph*),topsep=2pt,itemsep=1pt,partopsep=0pt,parsep=0pt]
    \item \relax \texttt{plt.bar()}
    \item \relax \texttt{plt.hist()}
    \item \relax \texttt{plt.barh()}
    \item \relax \texttt{plt.column()}
\end{enumerate}
\vspace{{2mm}}

\noindent\textbf{Q20.} \hfill {\small\color{LPUOrange}\textbf{[CO5 | pie charts]}}\\[2pt]
A pie chart is best suited for showing:

\begin{enumerate}[label=(\alph*),topsep=2pt,itemsep=1pt,partopsep=0pt,parsep=0pt]
    \item \relax Frequency distribution of bins
    \item \relax Proportions of a whole
    \item \relax Correlation between variables
    \item \relax Trends over time
\end{enumerate}
\vspace{{2mm}}

\noindent\textbf{Q21.} \hfill {\small\color{LPUOrange}\textbf{[CO6 | web-scraping-based url access]}}\\[2pt]
What does the term "web scraping" mean?

\begin{enumerate}[label=(\alph*),topsep=2pt,itemsep=1pt,partopsep=0pt,parsep=0pt]
    \item \relax Hosting websites on a server
    \item \relax Extracting data from websites automatically
    \item \relax Designing website layouts
    \item \relax Encrypting web traffic
\end{enumerate}
\vspace{{2mm}}

\noindent\textbf{Q22.} \hfill {\small\color{LPUOrange}\textbf{[CO6 | web-scraping-based url access]}}\\[2pt]
An HTTP status code of \texttt{200} typically indicates what?

\begin{enumerate}[label=(\alph*),topsep=2pt,itemsep=1pt,partopsep=0pt,parsep=0pt]
    \item \relax Access is forbidden
    \item \relax The request succeeded
    \item \relax A server error occurred
    \item \relax The page was not found
\end{enumerate}
\vspace{{2mm}}

\noindent\textbf{Q23.} \hfill {\small\color{LPUOrange}\textbf{[CO6 | data visualization]}}\\[2pt]
Which Python library is one of the most widely used for creating basic plots and charts?

\begin{enumerate}[label=(\alph*),topsep=2pt,itemsep=1pt,partopsep=0pt,parsep=0pt]
    \item \relax Matplotlib
    \item \relax Requests
    \item \relax PyGame
    \item \relax OpenCV
\end{enumerate}
\vspace{{2mm}}

\noindent\textbf{Q24.} \hfill {\small\color{LPUOrange}\textbf{[CO6 | data visualization]}}\\[2pt]
Which chart type is best suited for showing trends in data over time?

\begin{enumerate}[label=(\alph*),topsep=2pt,itemsep=1pt,partopsep=0pt,parsep=0pt]
    \item \relax Line chart
    \item \relax Pie chart
    \item \relax Word cloud
    \item \relax Scatter matrix
\end{enumerate}
\vspace{{2mm}}

\noindent\textbf{Q25.} \hfill {\small\color{LPUOrange}\textbf{[CO6 | data visualization]}}\\[2pt]
Which library is built on top of Matplotlib and provides attractive statistical graphics?

\begin{enumerate}[label=(\alph*),topsep=2pt,itemsep=1pt,partopsep=0pt,parsep=0pt]
    \item \relax BeautifulSoup
    \item \relax Seaborn
    \item \relax Flask
    \item \relax PyGame
\end{enumerate}
\vspace{{2mm}}

\noindent\textbf{Q26.} \hfill {\small\color{LPUOrange}\textbf{[CO6 | data visualization]}}\\[2pt]
Which chart is most appropriate for showing the proportion of parts to a whole?

\begin{enumerate}[label=(\alph*),topsep=2pt,itemsep=1pt,partopsep=0pt,parsep=0pt]
    \item \relax Histogram
    \item \relax Line chart
    \item \relax Pie chart
    \item \relax Box plot
\end{enumerate}
\vspace{{2mm}}

\noindent\textbf{Q27.} \hfill {\small\color{LPUOrange}\textbf{[CO6 | discussion on advanced toolboxes in python]}}\\[2pt]
Which library is commonly used for numerical computing and working with arrays in Python?

\begin{enumerate}[label=(\alph*),topsep=2pt,itemsep=1pt,partopsep=0pt,parsep=0pt]
    \item \relax Requests
    \item \relax PyGame
    \item \relax NumPy
    \item \relax Pillow
\end{enumerate}
\vspace{{2mm}}

\noindent\textbf{Q28.} \hfill {\small\color{LPUOrange}\textbf{[CO6 | discussion on advanced toolboxes in python]}}\\[2pt]
Which Python library is primarily used for data analysis and handling tabular data?

\begin{enumerate}[label=(\alph*),topsep=2pt,itemsep=1pt,partopsep=0pt,parsep=0pt]
    \item \relax OpenCV
    \item \relax Turtle
    \item \relax Pandas
    \item \relax PyGame
\end{enumerate}
\vspace{{2mm}}

\noindent\textbf{Q29.} \hfill {\small\color{LPUOrange}\textbf{[CO6 | demonstration of real-world applications]}}\\[2pt]
Which of the following is a real-world application where Python is commonly used?

\begin{enumerate}[label=(\alph*),topsep=2pt,itemsep=1pt,partopsep=0pt,parsep=0pt]
    \item \relax Manufacturing car engines
    \item \relax Data analysis and machine learning
    \item \relax Painting walls
    \item \relax Baking bread
\end{enumerate}
\vspace{{2mm}}

\noindent\textbf{Q30.} \hfill {\small\color{LPUOrange}\textbf{[CO6 | storytelling]}}\\[2pt]
In the context of data, what does "data storytelling" mainly involve?

\begin{enumerate}[label=(\alph*),topsep=2pt,itemsep=1pt,partopsep=0pt,parsep=0pt]
    \item \relax Encrypting datasets securely
    \item \relax Writing fictional novels with code
    \item \relax Deleting outdated records
    \item \relax Communicating insights using data and visuals
\end{enumerate}
\vspace{{2mm}}

\newpage
\subsection*{Part B: Descriptive Analytical \& Programming Section (40 Marks)}

\sffamily\textbf{Note: Attempt any FOUR questions from this section. Each question carries 10 marks.}
\vspace{{2mm}}

\questionitem{31}{Design a complete Python class \texttt{MatrixOperations} utilizing NumPy that performs matrix inversion, determinant calculation via LU decomposition, and solves a system of linear equations $Ax = b$ with error handling for singular matrices.}{10}{CO4}{L4 Analyze}
\vspace{3mm}

\questionitem{32}{Develop an end-to-end Pandas and Matplotlib data analytics pipeline to ingest a web transaction log, extract hourly traffic volume, compute 7-day moving averages, and plot a multi-curve trend with confidence intervals.}{10}{CO5}{L3 Apply}
\vspace{3mm}

\questionitem{33}{Design a real-time computer vision algorithm using OpenCV to track a colored object in a video stream, compute its bounding box centroid coordinates, and draw trajectory overlays.}{10}{CO6}{L3 Apply}
\vspace{3mm}

\questionitem{34}{Explain the complete Pygame event pump architecture. Construct an interactive game loop with a movable player sprite, wall collision boundaries, and frame-rate regulation at 60 FPS.}{10}{CO6}{L3 Apply}
\vspace{3mm}

\questionitem{35}{Design a resilient multi-threaded web scraper using \texttt{concurrent.futures}, \texttt{requests}, and \texttt{BeautifulSoup} to scrape research paper abstracts from an archive portal, handle HTTP timeouts, and persist structured JSON records.}{10}{CO6}{L4 Analyze}
\vspace{3mm}

\newpage
\subsection*{Official Solutions and Marking Rubrics --- End-Term Set B}

\subsubsection*{Part A Answer Key}

\begin{center}
\begin{tabularx}{0.96\textwidth}{c c X c c X}
\toprule
\textbf{Q\#} & \textbf{Ans} & \textbf{Core Topic} & \textbf{Q\#} & \textbf{Ans} & \textbf{Core Topic} \\
\midrule
1 & \textbf{B} & ndarray contiguous C-order buffer & 16 & \textbf{C} & Groupby split-apply-combine \\
2 & \textbf{A} & Vectorized broadcasting compatibility & 17 & \textbf{A} & Resample timeseries frequency \\
3 & \textbf{D} & Stride calculation for view & 18 & \textbf{D} & Merge join types inner/outer \\
4 & \textbf{C} & Fancy indexing produces copy & 19 & \textbf{B} & Matplotlib subplots OOP model \\
5 & \textbf{A} & Universal function ufunc & 20 & \textbf{C} & Histogram density normalization \\
6 & \textbf{B} & Reduction axis collapses dimension & 21 & \textbf{A} & OpenCV loads in BGR order \\
7 & \textbf{C} & Reshape total elements invariant & 22 & \textbf{B} & Color conversion COLOR\_BGR2GRAY \\
8 & \textbf{D} & Flatten copy vs ravel view & 23 & \textbf{D} & Gaussian blur 2D convolution \\
9 & \textbf{B} & Matrix product @ vs * Hadamard & 24 & \textbf{C} & Canny double threshold hysteresis \\
10 & \textbf{A} & Broadcasting dimensions alignment & 25 & \textbf{A} & Contour retrieval external \\
11 & \textbf{C} & Series 1D labeled array & 26 & \textbf{D} & Pygame canonical event loop \\
12 & \textbf{B} & DataFrame 2D tabular data & 27 & \textbf{B} & Clock tick locks FPS rate \\
13 & \textbf{A} & .loc inclusive vs .iloc exclusive & 28 & \textbf{C} & Requests HTTP status checking \\
14 & \textbf{D} & SettingWithCopyWarning prevention & 29 & \textbf{A} & BeautifulSoup DOM parser \\
15 & \textbf{B} & Missing value detection isna & 30 & \textbf{D} & Scikit-learn fit on train set only \\
\bottomrule
\end{tabularx}
\end{center}
\vspace{{4mm}}

\subsubsection*{Part B Comprehensive Solutions \& Marking Schemes}

\begin{solutionbox}[Question 31 Solution]
\begin{pycode}
import numpy as np\\
~\\
class MatrixOperations:\\
\_\_ESCTOKEN\_7\_\_ def \_\_ESCTOKEN\_0\_\_\_\_ESCTOKEN\_0\_\_init\_\_ESCTOKEN\_0\_\_\_\_ESCTOKEN\_0\_\_(self, A):\\
\_\_ESCTOKEN\_7\_\_ \_\_ESCTOKEN\_7\_\_ self.A = np.array(A, dtype=float)\\
\_\_ESCTOKEN\_7\_\_ \_\_ESCTOKEN\_7\_\_ if self.A.ndim != 2 or self.A.shape[0] != self.A.shape[1]:\\
\_\_ESCTOKEN\_7\_\_ \_\_ESCTOKEN\_7\_\_ \_\_ESCTOKEN\_7\_\_ raise ValueError("Matrix must be 2D square")\\
~\\
\_\_ESCTOKEN\_7\_\_ def determinant(self):\\
\_\_ESCTOKEN\_7\_\_ \_\_ESCTOKEN\_7\_\_ return np.linalg.det(self.A)\\
~\\
\_\_ESCTOKEN\_7\_\_ def solve\_\_ESCTOKEN\_0\_\_system(self, b):\\
\_\_ESCTOKEN\_7\_\_ \_\_ESCTOKEN\_7\_\_ b = np.array(b, dtype=float)\\
\_\_ESCTOKEN\_7\_\_ \_\_ESCTOKEN\_7\_\_ det = self.determinant()\\
\_\_ESCTOKEN\_7\_\_ \_\_ESCTOKEN\_7\_\_ if np.isclose(det, 0.0):\\
\_\_ESCTOKEN\_7\_\_ \_\_ESCTOKEN\_7\_\_ \_\_ESCTOKEN\_7\_\_ raise np.linalg.LinAlgError("Matrix is singular (det=0); no unique solution exists")\\
\_\_ESCTOKEN\_7\_\_ \_\_ESCTOKEN\_7\_\_ return np.linalg.solve(self.A, b)
\end{pycode}
\textbf{Complexity Analysis}: LU factorization requires $\mathcal{O}(n^3)$ floating-point operations. The \texttt{np.isclose} check prevents numerical instability artifacts in near-singular matrices.
\begin{markingrubric}
$\bullet$ Square matrix validation and OOP design: \textbf{3 Marks}\\$\bullet$ Determinant and LinAlgError exception handling: \textbf{4 Marks}\\$\bullet$ Asymptotic complexity $\mathcal{O}(n^3)$ derivation: \textbf{3 Marks}
\end{markingrubric}
\end{solutionbox}

\begin{solutionbox}[Question 32 Solution]
\begin{pycode}
import pandas as pd\\
import matplotlib.pyplot as plt\\
~\\
def analyze\_\_ESCTOKEN\_0\_\_traffic(csv\_\_ESCTOKEN\_0\_\_path):\\
\_\_ESCTOKEN\_7\_\_ df = pd.read\_\_ESCTOKEN\_0\_\_csv(csv\_\_ESCTOKEN\_0\_\_path, parse\_\_ESCTOKEN\_0\_\_dates=['timestamp'])\\
\_\_ESCTOKEN\_7\_\_ df.set\_\_ESCTOKEN\_0\_\_index('timestamp', inplace=True)\\
\_\_ESCTOKEN\_7\_\_ hourly = df.resample('H').size().rename('requests')\\
\_\_ESCTOKEN\_7\_\_ rolling\_\_ESCTOKEN\_0\_\_7d = hourly.rolling(24 * 7).mean()\\
\_\_ESCTOKEN\_7\_\_ std\_\_ESCTOKEN\_0\_\_7d = hourly.rolling(24 * 7).std()\\
~\\
\_\_ESCTOKEN\_7\_\_ plt.figure(figsize=(14, 6))\\
\_\_ESCTOKEN\_7\_\_ plt.plot(hourly.index, hourly, alpha=0.3, label='Hourly Traffic')\\
\_\_ESCTOKEN\_7\_\_ plt.plot(rolling\_\_ESCTOKEN\_0\_\_7d.index, rolling\_\_ESCTOKEN\_0\_\_7d, color='navy', label='7D Moving Average')\\
\_\_ESCTOKEN\_7\_\_ plt.fill\_\_ESCTOKEN\_0\_\_between(rolling\_\_ESCTOKEN\_0\_\_7d.index, rolling\_\_ESCTOKEN\_0\_\_7d - 2*std\_\_ESCTOKEN\_0\_\_7d, rolling\_\_ESCTOKEN\_0\_\_7d + 2*std\_\_ESCTOKEN\_0\_\_7d,\\
\_\_ESCTOKEN\_7\_\_ \_\_ESCTOKEN\_7\_\_ \_\_ESCTOKEN\_7\_\_ \_\_ESCTOKEN\_7\_\_ \_\_ESCTOKEN\_7\_\_ color='blue', alpha=0.15, label='95\_\_ESCTOKEN\_2\_\_ Confidence Band')\\
\_\_ESCTOKEN\_7\_\_ plt.title("Web Traffic Trends")\\
\_\_ESCTOKEN\_7\_\_ plt.legend(); plt.tight\_\_ESCTOKEN\_0\_\_layout()\\
\_\_ESCTOKEN\_7\_\_ plt.savefig("traffic.png")
\end{pycode}
\begin{markingrubric}
$\bullet$ Ingestion, timestamp parsing, and resampling: \textbf{3 Marks}\\$\bullet$ Rolling statistics computation: \textbf{3 Marks}\\$\bullet$ Multi-curve Matplotlib plot with shaded confidence region: \textbf{4 Marks}
\end{markingrubric}
\end{solutionbox}

\begin{solutionbox}[Question 33 Solution]
\begin{pycode}
import cv2\\
import numpy as np\\
~\\
def track\_\_ESCTOKEN\_0\_\_colored\_\_ESCTOKEN\_0\_\_object(frame):\\
\_\_ESCTOKEN\_7\_\_ hsv = cv2.cvtColor(frame, cv2.COLOR\_\_ESCTOKEN\_0\_\_BGR2HSV)\\
\_\_ESCTOKEN\_7\_\_ lower\_\_ESCTOKEN\_0\_\_blue = np.array([100, 150, 50])\\
\_\_ESCTOKEN\_7\_\_ upper\_\_ESCTOKEN\_0\_\_blue = np.array([140, 255, 255])\\
\_\_ESCTOKEN\_7\_\_ mask = cv2.inRange(hsv, lower\_\_ESCTOKEN\_0\_\_blue, upper\_\_ESCTOKEN\_0\_\_blue)\\
\_\_ESCTOKEN\_7\_\_ mask = cv2.erode(mask, None, iterations=2)\\
\_\_ESCTOKEN\_7\_\_ mask = cv2.dilate(mask, None, iterations=2)\\
~\\
\_\_ESCTOKEN\_7\_\_ contours, \_\_ESCTOKEN\_0\_\_ = cv2.findContours(mask, cv2.RETR\_\_ESCTOKEN\_0\_\_EXTERNAL, cv2.CHAIN\_\_ESCTOKEN\_0\_\_APPROX\_\_ESCTOKEN\_0\_\_SIMPLE)\\
\_\_ESCTOKEN\_7\_\_ if contours:\\
\_\_ESCTOKEN\_7\_\_ \_\_ESCTOKEN\_7\_\_ c = max(contours, key=cv2.contourArea)\\
\_\_ESCTOKEN\_7\_\_ \_\_ESCTOKEN\_7\_\_ ((x, y), radius) = cv2.minEnclosingCircle(c)\\
\_\_ESCTOKEN\_7\_\_ \_\_ESCTOKEN\_7\_\_ M = cv2.moments(c)\\
\_\_ESCTOKEN\_7\_\_ \_\_ESCTOKEN\_7\_\_ if M["m00"] > 0:\\
\_\_ESCTOKEN\_7\_\_ \_\_ESCTOKEN\_7\_\_ \_\_ESCTOKEN\_7\_\_ cx = int(M["m10"] / M["m00"])\\
\_\_ESCTOKEN\_7\_\_ \_\_ESCTOKEN\_7\_\_ \_\_ESCTOKEN\_7\_\_ cy = int(M["m01"] / M["m00"])\\
\_\_ESCTOKEN\_7\_\_ \_\_ESCTOKEN\_7\_\_ \_\_ESCTOKEN\_7\_\_ cv2.circle(frame, (cx, cy), 5, (0, 0, 255), -1)\\
\_\_ESCTOKEN\_7\_\_ return frame
\end{pycode}
\begin{markingrubric}
$\bullet$ Color space conversion to HSV and thresholding: \textbf{3 Marks}\\$\bullet$ Morphological filtering (erosion/dilation): \textbf{3 Marks}\\$\bullet$ Centroid calculation via spatial moments ($m_{10}/m_{00}$): \textbf{4 Marks}
\end{markingrubric}
\end{solutionbox}

\begin{solutionbox}[Question 34 Solution]
\begin{pycode}
import pygame, sys\\
~\\
pygame.init()\\
screen = pygame.display.set\_\_ESCTOKEN\_0\_\_mode((800, 600))\\
clock = pygame.time.Clock()\\
player = pygame.Rect(380, 280, 40, 40)\\
speed = 5\\
~\\
while True:\\
\_\_ESCTOKEN\_7\_\_ for event in pygame.event.get():\\
\_\_ESCTOKEN\_7\_\_ \_\_ESCTOKEN\_7\_\_ if event.type == pygame.QUIT:\\
\_\_ESCTOKEN\_7\_\_ \_\_ESCTOKEN\_7\_\_ \_\_ESCTOKEN\_7\_\_ pygame.quit(); sys.exit()\\
~\\
\_\_ESCTOKEN\_7\_\_ keys = pygame.key.get\_\_ESCTOKEN\_0\_\_pressed()\\
\_\_ESCTOKEN\_7\_\_ if keys[pygame.K\_\_ESCTOKEN\_0\_\_LEFT]:  player.x -= speed\\
\_\_ESCTOKEN\_7\_\_ if keys[pygame.K\_\_ESCTOKEN\_0\_\_RIGHT]: player.x += speed\\
\_\_ESCTOKEN\_7\_\_ if keys[pygame.K\_\_ESCTOKEN\_0\_\_UP]:    player.y -= speed\\
\_\_ESCTOKEN\_7\_\_ if keys[pygame.K\_\_ESCTOKEN\_0\_\_DOWN]:  player.y += speed\\
~\\
\_\_ESCTOKEN\_7\_\_ \_\_ESCTOKEN\_1\_\_ Boundary clipping:\\
\_\_ESCTOKEN\_7\_\_ player.clamp\_\_ESCTOKEN\_0\_\_ip(screen.get\_\_ESCTOKEN\_0\_\_rect())\\
~\\
\_\_ESCTOKEN\_7\_\_ screen.fill((30, 30, 30))\\
\_\_ESCTOKEN\_7\_\_ pygame.draw.rect(screen, (0, 200, 255), player)\\
\_\_ESCTOKEN\_7\_\_ pygame.display.flip()\\
\_\_ESCTOKEN\_7\_\_ clock.tick(60)
\end{pycode}
\begin{markingrubric}
$\bullet$ Pygame initialization and surface management: \textbf{2 Marks}\\$\bullet$ Event pump handling and graceful shutdown: \textbf{3 Marks}\\$\bullet$ Boundary collision clamp and 60 FPS clock regulator: \textbf{5 Marks}
\end{markingrubric}
\end{solutionbox}

\begin{solutionbox}[Question 35 Solution]
\begin{pycode}
import requests, json\\
from bs4 import BeautifulSoup\\
from concurrent.futures import ThreadPoolExecutor\\
~\\
def fetch\_\_ESCTOKEN\_0\_\_abstract(paper\_\_ESCTOKEN\_0\_\_id):\\
\_\_ESCTOKEN\_7\_\_ url = f"https://example-archive.org/paper/\_\_ESCTOKEN\_5\_\_paper\_\_ESCTOKEN\_0\_\_id\_\_ESCTOKEN\_6\_\_"\\
\_\_ESCTOKEN\_7\_\_ try:\\
\_\_ESCTOKEN\_7\_\_ \_\_ESCTOKEN\_7\_\_ resp = requests.get(url, timeout=5, headers=\_\_ESCTOKEN\_5\_\_'User-Agent': 'Bot/1.0'\_\_ESCTOKEN\_6\_\_)\\
\_\_ESCTOKEN\_7\_\_ \_\_ESCTOKEN\_7\_\_ if resp.status\_\_ESCTOKEN\_0\_\_code == 200:\\
\_\_ESCTOKEN\_7\_\_ \_\_ESCTOKEN\_7\_\_ \_\_ESCTOKEN\_7\_\_ soup = BeautifulSoup(resp.text, 'html.parser')\\
\_\_ESCTOKEN\_7\_\_ \_\_ESCTOKEN\_7\_\_ \_\_ESCTOKEN\_7\_\_ title = soup.find('h1', class\_\_ESCTOKEN\_0\_\_='title').get\_\_ESCTOKEN\_0\_\_text(strip=True)\\
\_\_ESCTOKEN\_7\_\_ \_\_ESCTOKEN\_7\_\_ \_\_ESCTOKEN\_7\_\_ abstract = soup.find('div', class\_\_ESCTOKEN\_0\_\_='abstract').get\_\_ESCTOKEN\_0\_\_text(strip=True)\\
\_\_ESCTOKEN\_7\_\_ \_\_ESCTOKEN\_7\_\_ \_\_ESCTOKEN\_7\_\_ return \_\_ESCTOKEN\_5\_\_'id': paper\_\_ESCTOKEN\_0\_\_id, 'title': title, 'abstract': abstract\_\_ESCTOKEN\_6\_\_\\
\_\_ESCTOKEN\_7\_\_ except requests.exceptions.RequestException:\\
\_\_ESCTOKEN\_7\_\_ \_\_ESCTOKEN\_7\_\_ return None\\
~\\
def run\_\_ESCTOKEN\_0\_\_scraper(paper\_\_ESCTOKEN\_0\_\_ids):\\
\_\_ESCTOKEN\_7\_\_ with ThreadPoolExecutor(max\_\_ESCTOKEN\_0\_\_workers=5) as executor:\\
\_\_ESCTOKEN\_7\_\_ \_\_ESCTOKEN\_7\_\_ results = list(executor.map(fetch\_\_ESCTOKEN\_0\_\_abstract, paper\_\_ESCTOKEN\_0\_\_ids))\\
\_\_ESCTOKEN\_7\_\_ valid = [r for r in results if r is not None]\\
\_\_ESCTOKEN\_7\_\_ with open('papers.json', 'w') as f:\\
\_\_ESCTOKEN\_7\_\_ \_\_ESCTOKEN\_7\_\_ json.dump(valid, f, indent=2)
\end{pycode}
\begin{markingrubric}
$\bullet$ Resilient HTTP requests with timeouts and exception catching: \textbf{3 Marks}\\$\bullet$ Concurrent execution via ThreadPoolExecutor: \textbf{4 Marks}\\$\bullet$ DOM parsing and structured JSON persistence: \textbf{3 Marks}
\end{markingrubric}
\end{solutionbox}


\end{document}
